% Capítulo 52 del sucesor: integración pública del propietario íntegro del radión.
% No carga portada, main, style ni C_estatutos_procedencia del módulo autónomo.

% Adaptador de compatibilidad del propietario íntegro del radión.
%
% El tratado rector ya carga tipografía, geometría, hipervínculos, teoremas,
% tablas, TikZ y tcolorbox.  Este archivo aporta únicamente los nombres y
% entornos exclusivos del propietario, sin importar su preámbulo autónomo ni
% alterar la jerarquía tipográfica general.

\makeatletter
\@ifundefined{color@RadNavy}{\definecolor{RadNavy}{HTML}{102A43}}{}
\@ifundefined{color@RadBlue}{\definecolor{RadBlue}{HTML}{1F5E7A}}{}
\@ifundefined{color@RadTeal}{\definecolor{RadTeal}{HTML}{168C8C}}{}
\@ifundefined{color@RadCopper}{\definecolor{RadCopper}{HTML}{B66A3C}}{}
\@ifundefined{color@RadGold}{\definecolor{RadGold}{HTML}{D4A73A}}{}
\@ifundefined{color@RadInk}{\definecolor{RadInk}{HTML}{1A202C}}{}
\@ifundefined{color@RadMist}{\definecolor{RadMist}{HTML}{EDF4F7}}{}
\@ifundefined{color@RadCream}{\definecolor{RadCream}{HTML}{FAF7F0}}{}
\@ifundefined{color@RadRose}{\definecolor{RadRose}{HTML}{8D3F55}}{}
\@ifundefined{color@RadGreen}{\definecolor{RadGreen}{HTML}{2F7D5D}}{}
\@ifundefined{color@RadGray}{\definecolor{RadGray}{HTML}{66788A}}{}

\providecommand{\TRIT}{\ensuremath{\mathrm{TRIT}}}
\providecommand{\HMT}{\ensuremath{\mathrm{HMT}}}
\providecommand{\Rstate}{\mathcal R_{12}}
\providecommand{\epsR}{\varepsilon_R}
\providecommand{\epsone}{\varepsilon_R^{(1)}}
\providecommand{\pih}{\pi_{\!\HMT}}
\providecommand{\alphah}{\alpha_{\!\HMT}}
\providecommand{\hret}{\hbar_{\mathrm{ret}}}
\providecommand{\hfull}{h_{\mathrm{ret}}}
\providecommand{\diff}{\mathop{}\!\mathrm d}
\providecommand{\e}{\mathrm e}
\providecommand{\R}{\mathbb R}
\providecommand{\C}{\mathbb C}
\providecommand{\F}{\mathbb F}
\providecommand{\Z}{\mathbb Z}
\providecommand{\dd}{\,\mathrm d}
\providecommand{\status}[1]{\textsf{\bfseries #1}}
\providecommand{\sourcepath}[1]{\nolinkurl{#1}}

\@ifundefined{typebox}{%
  \newtcolorbox{typebox}[1][]{enhanced,breakable,colback=white,
    colframe=RadBlue,boxrule=.7pt,arc=1.5mm,left=4mm,right=4mm,
    top=3mm,bottom=3mm,title={Reçu de types},fonttitle=\bfseries,
    coltitle=RadNavy,#1}%
}{}
\@ifundefined{narrativebox}{%
  \newtcolorbox{narrativebox}[1][]{enhanced,breakable,colback=white,
    colframe=RadGold,borderline west={2.4pt}{0pt}{RadGold},
    boxrule=.25pt,arc=0mm,left=5mm,right=4mm,top=3mm,bottom=3mm,#1}%
}{}

% Las once portadillas internas del propietario autónomo se convierten en
% divisiones científicas ordinarias dentro del capítulo 52.  Se evita así
% introducir partes autónomas, gigantismo tipográfico o páginas vacías.
\providecommand{\partopener}[3]{\section{#2}\noindent\emph{#3}\par\medskip}

\providecommand{\BigRadionEquation}{%
  \begin{equation*}
  \boxed{\displaystyle
  \frac{180^{\circ}}{\pih}
  =50^{\circ}+A^{\circ}-\frac{20}{81}\,\Delta_4^{\circ}+\epsR^{\circ}}
  \end{equation*}}
\makeatother


\chapter{Le radion angulaire HMT--MD : raccord avec retenue des sections directe et torsionnelle et clôture exacte de l'unité de phase, d'action et de mémoire}
\label{chap:sucesor-102-radion}

\section{Composition APP--TRIT--TPK de l'état de phase : génération de l'unité orientée avec mémoire}
% Fragmento público generado desde propietarios íntegros.
% Residencia: 52.1.
% Fuente: ../incoming/radion_monografia_20260827/manuscrito/sections/01_resumen_y_guia.tex:1-128
\paragraph{Synthèse mathématique de l'objet et de sa chaîne causale}

Cette monographie reconstruit le radian à partir de la chaîne générative
\[
\APP\longrightarrow\TRIT\longrightarrow\TPK
\longrightarrow X_{\mathrm{enr}}
\longrightarrow\text{structure discrète du continu}
\longrightarrow\MD,
\]
sans prendre comme entrée ni la valeur métrologique de \(180/\pi\), ni le petit
résidu angulaire que l'on souhaite expliquer. Le résultat est une unité enrichie,
le \emph{radion}, défini comme l'unité orientée qui parcourt un radian de
phase et balaie exactement une action \(\hret\) dans l'ellipse positive associée
au même état.

La clôture directrice est
\BigRadionEquation
avec
\[
S_E=2\pih\left(\frac5{36}+\frac{25}{9}\alphah\right),
\qquad
\Delta_4=
\frac{\pih-e_{\HMT}\log\pih}{270}
\left(1+\frac{\pih}{729}\right),
\]
\[
\rho_{\mathrm{tw}}=\frac{2\cdot90}{729}=\frac{20}{81},
\qquad
\epsone=\frac{20}{81}\Delta_4-(S_E-1).
\]
La division APP prouve \(1<S_E<2\), de sorte que le quotient est exactement
un et le reste est \(D_E=S_E-1\). L'égalité
\[
1=S_E-\frac{20}{81}\Delta_4+\epsone
\]
est alors une identité interne. La carte angulaire ultérieure produit
\[
\epsR^\circ
=5.13830167585752820716851646226459474899085744\ldots
\times10^{-8}\,{}^\circ.
\]

Le texte démontre en outre que l'ellipse d'action
\[
Q=a_S\cos\theta,\qquad P=b_S\sin\theta,\qquad
a_Sb_S=2\hret
\]
satisfait
\[
\operatorname{Area}(0,\theta)=\hret\theta.
\]
Par conséquent, un radion balaie \(\hret\) et le tour \(2\pi\) enferme
\(h_{\mathrm{ret}}=2\pi\hret\). La même ellipse possède un bord symplectique
fini et un intérieur d'aire infinie lorsqu'elle est munie de la métrique
de Hilbert--Beltrami--Klein. Telle est la forme exacte de la thèse physique qui guide
l'ouvrage : la surface visible clôt l'action ; la profondeur intérieure
conserve un développement non épuisable par la projection plane.

L'exposition maintient typées quatre réalisations du même état :

\begin{enumerate}[label=\textbf{\arabic*.}]
\item le cercle \(S^1\), qui publie la phase ;
\item l'ellipse positive, qui publie l'action et l'anisotropie ;
\item l'intérieur de Klein, qui mesure la profondeur au moyen du birapport ;
\item l'hélice nonadique, qui conserve retour, mémoire et contraction.
\end{enumerate}

Ce ne sont pas quatre géométries superposées sans critère. Ce sont quatre lecteurs ayant
des domaines, des codomaines et des données conservées différents, appliqués à un unique état
enrichi. Le secteur \(90/120\), la dodécaphase, la retenue du radion, la queue
de Lambert, le noyau de Barbero et les tours globales sont également présentés
comme des objets apparentés mais non interchangeables.

\begin{thesisbox}
Le radion est l'unité orientée qui balaie \(\hret\) dans l'ellipse d'action.
Le cercle publie sa phase, Klein mesure sa profondeur, l'hélice conserve sa
mémoire et la retenue coinductive raccorde exactement ses sections directe et
torsionnelle.
\end{thesisbox}

\paragraph{Ordre d'exposition et dépendances de lecture}

L'ordre n'est pas ornemental. Les Parties I et II produisent l'objet et sa clôture
à partir d'APP--TRIT--TPK. Ce n'est qu'ensuite que sont introduites les réalisations dans le langage
de Hilbert, LC, Maxwell, Minkowski, Hodge, Lorentz, de la modularité et des orbites. Chaque
chapitre distingue trois niveaux :

\begin{description}[style=nextline,leftmargin=4.2cm,labelwidth=3.9cm]
\item[Résultat interne.] Égalité ou construction démontrée dans le domaine
HMT avec ses primitives déclarées.
\item[Formalisation réunie.] Nouvelle application articulant des résultats déjà existants
sans utiliser la cible comme sélecteur.
\item[Interface physique.] Traduction dimensionnelle ultérieure reconnaissant une
structure conventionnelle sans en faire l'origine de la sélection.
\end{description}

Les encadrés de \emph{statut et falsificateur} indiquent quelle observation concrète
invaliderait chaque promotion forte. Cette discipline n'affaiblit pas la thèse ; elle la
rend auditable. L'objectif de la narration étendue est qu'aucun mot
comme \emph{cercle}, \emph{ellipse}, \emph{torsion}, \emph{charge} ou
\emph{mémoire} ne masque un changement de type.

\paragraph{Résultats directeurs et portée probatoire}

\begin{enumerate}[label=\textbf{C\arabic*.}]
\item Clôture indépendante de la valeur finale du radion par deux publications du même état et
soustraction avec retenue APP.
\item Dérivation du coefficient \(20/81\) comme deux feuillets du secteur actif
\(90\), normalisés par la fibre \(3^6=729\).
\item Égalisateur typé entre la couture \(\Re s=1/2\), la face APP de cent
marques et la phase dodécaphasique \(\theta_2=50^\circ\).
\item Théorème d'aire du radion : \(\operatorname{Area}(\theta)=\hret\theta\).
\item Étude focale exacte de l'ellipse et de ses invariants euclidiens,
hyperboliques, symplectiques et modulaires.
\item Théorème bord--intérieur : action totale finie \(h\) et intérieur de Klein
d'aire hyperbolique infinie.
\item Réalisation exacte par covariance minimale et oscillateur LC, suivie de
la publication électromagnétique \(90/120\).
\item Atlas de non-identifications pour tous les objets \(90/120\), y compris
la séparation entre \(\epsR\), queue spectrale et noyau de Barbero \(12:-1\).
\item Caractérisation structurelle de \(\pi\) dans la carte du radion, distincte
du générateur coinductif primaire.
\item Synthèse ontologique : dimension comme rang minimal de mémoires qu'une
projection de phase ne peut plus conserver.
\end{enumerate}
% Fuente: ../incoming/radion_monografia_20260827/manuscrito/sections/01b_mapa_resultados.tex:1-133
\paragraph{Correspondance entre résultats, types et résidences causales}

La correspondance suivante ordonne les affirmations directrices et leurs démonstrations,
en maintenant visible la hiérarchie qui empêche que
la richesse physique efface la clôture mathématique qui la soutient.

% HMT_RESULT_BEGIN:RDN-01-GENEALOGIA:theorem
\begin{exactbox}[title={Généalogie APP--TRIT--TPK du radion}]
Le radion naît de la composition effective
\(\APP\to\TRIT\to\TPK\to X_{\mathrm{enr}}\to\MD\) ; il ne s'obtient pas
en redéfinissant rétrospectivement le radian conventionnel.
\end{exactbox}
% HMT_RESULT_END:RDN-01-GENEALOGIA:theorem

% HMT_RESULT_BEGIN:RDN-02-CIERRE:theorem
\begin{exactbox}[title={Clôture exacte et indépendance prospective}]
La retenue APP produit intérieurement
\[
1=S_E-\frac{20}{81}\Delta_4+\varepsilon_R^{(1)},
\qquad
\varepsilon_R^{(1)}=\frac{20}{81}\Delta_4-(S_E-1),
\]
et la carte angulaire transporte cette même identité vers
\(180^\circ/\pi_{\!\HMT}\) sans introduire le résidu cible.
\end{exactbox}
% HMT_RESULT_END:RDN-02-CIERRE:theorem

% HMT_RESULT_BEGIN:RDN-03-BISAGRA-50:theorem
\begin{exactbox}[title={Coordonnée critique de \(50^\circ\)}]
La coordonnée critique \(\Re s=1/2\), publiée sur la face APP de cent
marques et sur la roue dodécaphasique, sélectionne la phase
\(\theta_2=50^\circ\). C'est une égalité de lecteurs typés du même état.
\end{exactbox}
% HMT_RESULT_END:RDN-03-BISAGRA-50:theorem

% HMT_RESULT_BEGIN:RDN-04-COEFICIENTE-2081:theorem
\begin{exactbox}[title={Densité bidirectionnelle nonadique \(20/81\)}]
Le coefficient torsionnel n'est pas ajusté sur le réseau \(1/81\) : il est construit comme
\(\rho_{\mathrm{tw}}=2N_6/3^6=2\cdot90/729=20/81\), en conservant feuillet,
orientation et normalisation de la fibre nonadique.
\end{exactbox}
% HMT_RESULT_END:RDN-04-COEFICIENTE-2081:theorem

% HMT_RESULT_BEGIN:RDN-05-AREA-ACCION:theorem
\begin{exactbox}[title={Loi surfacique du radion}]
Pour l'ellipse positive
\(Q=a_S\cos\theta\), \(P=b_S\sin\theta\),
\(a_Sb_S=2\hbar_{\mathrm{ret}}\), l'aire orientée balayée est
\(\mathcal A(\theta)=\hbar_{\mathrm{ret}}\theta\). Un radion balaie exactement
\(\hbar_{\mathrm{ret}}\), et le tour \(2\pi\) enferme
\(h_{\mathrm{ret}}=2\pi\hbar_{\mathrm{ret}}\).
\end{exactbox}
% HMT_RESULT_END:RDN-05-AREA-ACCION:theorem

% HMT_RESULT_BEGIN:RDN-06-INVARIANTES-FOCALES:theorem
\begin{exactbox}[title={Invariants focaux de l'opérateur positif}]
Le couple publié \(A\pm C^\ast\) fixe les demi-axes, l'excentricité, la
séparation focale et le paramètre de compression symplectique. En particulier,
\((d_1/d_0)^2=C^\ast/2\), identité adimensionnelle qui sépare la forme de l'échelle.
\end{exactbox}
% HMT_RESULT_END:RDN-06-INVARIANTES-FOCALES:theorem

% HMT_RESULT_BEGIN:RDN-07-BORDE-INTERIOR:theorem
\begin{exactbox}[title={Finitude symplectique du bord et divergence de l'aire hyperbolique intérieure}]
Le bord symplectique de l'ellipse enferme l'action finie
\(h_{\mathrm{ret}}\), tandis que l'aire de Hilbert--Beltrami--Klein de l'intérieur
diverge lorsqu'elle atteint la frontière. Phase visible et profondeur ne sont pas la même
mesure.
\end{exactbox}
% HMT_RESULT_END:RDN-07-BORDE-INTERIOR:theorem

% HMT_RESULT_BEGIN:RDN-08-HELICE-MEMORIA:theorem
\begin{exactbox}[title={Relèvement hélicoïdal nonadique avec retour de phase et avance de mémoire}]
La monodromie nonadique ramène la phase sans réinitialiser l'état. L'hélice
relève le cercle en conservant retenue et mémoire ; la clôture \(2\pi/4\pi\)
distingue retour projectif, inversion d'orientation et retour complet.
\end{exactbox}
% HMT_RESULT_END:RDN-08-HELICE-MEMORIA:theorem

% HMT_RESULT_BEGIN:RDN-09-OSCILADOR-EM:development
\begin{narrativebox}[title={Réalisation oscillatoire et constitutive électromagnétique}]
La même anisotropie se réalise par une covariance minimale et un oscillateur
LC : une coordonnée stocke l'aspect électrique-élastique et sa conjuguée
l'aspect magnétique-rotationnel. La fréquence et l'impédance sont typées
séparément.
\end{narrativebox}
% HMT_RESULT_END:RDN-09-OSCILADOR-EM:development

% HMT_RESULT_BEGIN:RDN-10-TIPOS-90120:theorem
\begin{exactbox}[title={Séparation typée des opérateurs \(90/120\)}]
Extracteur dodécaphasique, retenue du radion, queue de Lambert, fonctionnelle
hexade--octade, noyau de Barbero et semi-groupe global de tours sont des objets
distincts. Le poids \(12:-1\) procède des douze secteurs de la chaîne
fondamentale dodécaphasique et de son unique couture orientée. Le coefficient de
\((1-q)^{-1}\), égal à \(1/24\), et le contre-terme entier minimal vérifient
ultérieurement le couple produit.
\end{exactbox}
% HMT_RESULT_END:RDN-10-TIPOS-90120:theorem

% HMT_RESULT_BEGIN:RDN-11-GEOMETRIAS-CAMPO:development
\begin{narrativebox}[title={Réalisations complexe, lorentzienne et duale de Hodge}]
Les régimes elliptique, parabolique et hyperbolique du TRIT sont publiés comme
rotation, seuil et transformation hyperbolique. L'étoile de Hodge, les lois constitutives et la force
de Lorentz agissent ensuite sur l'état déjà orienté ; elles ne sélectionnent pas ses
constantes.
\end{narrativebox}
% HMT_RESULT_END:RDN-11-GEOMETRIAS-CAMPO:development

% HMT_RESULT_BEGIN:RDN-12-MODULAR-ORBITAL:development
\begin{narrativebox}[title={Caractères modulaire, orbital et arithmétique de frontière}]
L'involution modulaire, les orbites elliptiques, l'avance-retard et l'arithmétique
3--6--9 reconnaissent différentes mémoires du même état. La fonction modulaire et
les vacances sont utilisées comme lecteurs ultérieurs, jamais comme substituts du
générateur APP--TRIT--TPK.
\end{narrativebox}
% HMT_RESULT_END:RDN-12-MODULAR-ORBITAL:development

% HMT_RESULT_BEGIN:RDN-13-SINTESIS-ONTOLOGICA:conclusion
\begin{thesisbox}[title={Synthèse opératoire de phase, d'action, de profondeur et de mémoire}]
Le cercle est la projection de phase ; l'ellipse est la section d'action ; Klein
mesure la profondeur ; l'hélice conserve l'histoire. Le radion est l'unité
orientée qui maintient ces publications couplées sans effacer le résidu qui
les distingue.
\end{thesisbox}
% HMT_RESULT_END:RDN-13-SINTESIS-ONTOLOGICA:conclusion

% HMT_RESULT_BEGIN:RDN-14-CONTROL-ALCANCE:control
\begin{statusbox}[title={Délimitation des domaines et continuité causale}]
Les lois complètes des masses, le raffinement électronique, CKM/CP et le
couloir exceptionnel conservent leurs propres démonstrations directrices. Cette monographie les
cite seulement là où ils partagent un antécédent ou une carte ; elle ne transfère pas au radion leurs
sélecteurs spécifiques et ne confond pas leurs résidences probatoires.
\end{statusbox}
% HMT_RESULT_END:RDN-14-CONTROL-ALCANCE:control


\subsection{Support bivarié d'APP, orientation ternaire du TRIT et actualisation de l'état enrichi par le TPK}
% Fragmento público generado desde propietarios íntegros.
% Residencia: 52.1.1.
% Fuente: ../incoming/radion_monografia_20260827/manuscrito/sections/02_genealogia.tex:1-180
\noindent La généalogie opératoire précède la mesure angulaire et détermine les variables conservées par chaque lecteur.\par\medskip

\subsubsection{Support APP--TRIT--TPK et construction de l'état enrichi}

\paragraph{Antériorité de l'état enrichi par rapport à la publication circulaire}

La définition scolaire du radian commence par un cercle déjà donné, un centre déjà
donné et une longueur d'arc déjà mesurable. Si le rayon et l'arc ont la
même longueur, l'angle sous-tendu mesure un radian. La définition est correcte
au sein de la géométrie euclidienne, mais ne répond pas à une question antérieure :
quelle structure permet à un tour de posséder phase, orientation, action et
mémoire, et pourquoi l'unité naturelle est-elle liée à \(2\pi\) ?

HMT inverse l'ordre causal. Elle construit d'abord l'état fini, puis sa
prolongation coinductive, ensuite le tour et seulement à la fin sa publication
angulaire. Le radion ne corrige pas le radian conventionnel ; il le relève vers son domaine
généalogique.

\begin{typebox}
\[
\APP\to\TRIT\to\TPK\to X_{\mathrm{enr}}
\to\mathcal C^{\mathrm{disc}}_{\mathrm{cont}}
\to\bigl(\pih,e_{\HMT},\varphi_{\HMT},\alphah,A,C^*,\Delta_4,\hret\bigr)
\to\mathfrak r_\sigma.
\]
Les mathématiques conventionnelles n'apparaissent qu'ensuite comme langage de réalisation :
\(S^1\), ellipse symplectique, métrique de Klein, espace de Hilbert, circuit
LC, écran de Minkowski et équations de champ.
\end{typebox}

\paragraph{Support bivarié d'APP : feuillets, résidus, quotients et retenue}

La plateforme d'arithmétique positionnelle (APP) conserve simultanément deux
réalisations des symboles \(1,\dots,9\) : un feuillet additif et un feuillet
multiplicatif. Un nombre n'est pas réduit à sa valeur publiée. Avec la valeur
subsistent feuillet, orientation, quotient, reste, retenue, frontière, chemin et
survie.

Le recensement de base comporte trois strates qui ne doivent jamais être identifiées par simple
coïncidence de cardinalité :
\[
104,976\longrightarrow468\;U_6\longrightarrow243\;w_6
\subset\F_3^6,\qquad |\F_3^6|=729.
\]
Les \(468\) émissions décimales, les \(243\) mots visibles et l'espace ambiant
de \(729\) mots sont des objets distincts. L'espace ambiant sera décisif pour la
normalisation \(20/81\).

En base cubique \(B=1000\), la soustraction orientée de deux mots de triades
s'effectue de droite à gauche :
\begin{equation}
u_i=t_i-v_i+c_{i+1},\qquad
r_i=u_i\bmod1000,\qquad
c_i=\frac{u_i-r_i}{1000}.
\label{eq:subacarreo}
\end{equation}
Les opérandes et la retenue distante étant fixés, quotient et reste sont uniques. Tel est
l'opérateur qui, plus loin, publiera les chiffres de \(\epsone\). Il n'y a
aucune décimale choisie après avoir vu le radian.

La complétion
\[
10^3=729+243+27+1
\]
n'est pas une curiosité décimale. Elle exprime la conservation de l'unité lors du passage
de l'espace ambiant ternaire à une chambre de mille positions. La différence
\(1000-729=271\) est la couronne complétante ; l'identité correcte est
\(999=729+270\), non \(729+271\).

\paragraph{Orientation ternaire et classification des 3 régimes opératoires du TRIT}

Le TRIT attribue à chaque état local l'un de trois régimes orientés. Dans la
convention temporelle active,
\[
\begin{aligned}
+1&\longleftrightarrow\text{retour, mémoire, passé},\\
0&\longleftrightarrow\text{seuil, présent},\\
-1&\longleftrightarrow\text{ouverture bidirectionnelle, futur}.
\end{aligned}
\]
La réalisation quadratique distingue :
\begin{equation}
J^2=-I,\qquad N^2=0,\qquad R^2=I.
\label{eq:tres-regimenes}
\end{equation}
Ainsi apparaissent, respectivement, la rotation elliptique, le seuil parabolique et
la séparation hyperbolique. On ne mélange pas trois métriques. On conserve trois
actions typées sur un antécédent commun.

La racine interne de \(-1\) n'est pas importée en supposant d'avance le plan
complexe. Dans le bloc APP, deux projections \(P,Q\) de rapport
\(r=3/4=90/120\) produisent
\[
J_{\mathrm{comm}}=\frac{[P,Q]}{\sqrt{r(1-r)}},\qquad J_{\mathrm{comm}}^2=-I,
\]
tandis que
\[
R_{\mathrm{comm}}=2P-I,\qquad R_{\mathrm{comm}}^2=I,\qquad
R_{\mathrm{comm}}J_{\mathrm{comm}}=-J_{\mathrm{comm}}R_{\mathrm{comm}}.
\]
Le couple génère la forme minimale de \(\mathrm{Cl}_{1,1}\simeq M_2(\R)\) :
rotation et ouverture émergent ensemble, mais restent des opérations différentes.

\paragraph{Composition dynamique du TPK : sélection, transport, mémoire et retour}

Le Topological Phase Kernel n'est pas le nom d'une boîte d'objets. C'est la
composition qui sélectionne la phase, transporte la cardinalité, relève les feuillets,
enregistre la mémoire et exécute le retour. La dépendance fondamentale est
\[
U_t\longrightarrow\Gamma_9\longrightarrow K_{\mathrm{ph}}.
\]
L'holonomie nonadique \(\Gamma_9\) agit sur l'état enrichi ; la
publication \(K_{\mathrm{ph}}\) est ultérieure et conserve une mémoire réduite. L'identité
\(K_{\mathrm{ph}}^9|_r=E_+E_\times\) ne transforme pas \(K_{\mathrm{ph}}\) en
générateur.

Le certificat nonadique conserve la chaîne complète
\[
w_6\to w_{12}\to w_{18}\to w_{24}\to w_{30}
\to R_{36}\to G_9,
\]
avec \(396\) niveaux, \(2376\) trits, \(43\) tours, \(387\) paires
\((K,K+9)\) par canal et cinq régions de \(\pi\). La phase satisfait
\[
g(k+9)=g(k),
\]
mais l'état complet ne revient pas :
\[
B_{k+9}\ne B_k.
\]
Tel est le germe mathématique de l'hélice : la projection angulaire se ferme ; la
mémoire avance.

\paragraph{Extension temporelle non scindée \(C_{12}\to C_{108}\to C_9\)}

La mémoire de douze pas s'organise au moyen de
\begin{equation}
0\longrightarrow C_{12}\longrightarrow C_{108}
\longrightarrow C_9\longrightarrow0.
\label{eq:extension-108}
\end{equation}
Le cocycle
\[
c(r,r')=\left[\left\lfloor\frac{r+r'}9\right\rfloor\right]_{12}
\]
mesure la retenue entre retour de phase et avance de mémoire. Neuf pas ferment
la phase observable \(C_9\) ; ils ne réinitialisent pas le registre dodécaphasique.

Il convient de séparer quatre objets :
\begin{center}
\begin{tabularx}{.95\textwidth}{>{\bfseries}l X}
\toprule
\(C_{12}\) & groupe cyclique de douze positions ;\\
\(\Rstate\) & registre enrichi \((E_m,\tau_m,\sigma_m,\kappa_m,\Lambda_m)_{m=1}^{12}\) ;\\
\(\Theta_{108}\) & publication brute dans le calendrier de 108 ;\\
\(\Gamma_{108}\) & histoire ou holonomie relevée.\\
\bottomrule
\end{tabularx}
\end{center}
C'est pourquoi l'expression vague \enquote{post-\(C_{12}\)} sera remplacée par
\enquote{postérieur au registre dodécaphasique \(\Rstate\), dans sa lecture
sexadique} lorsque telle sera l'application effective.

\paragraph{Émergence corrélative des 5 constructions du continu à partir de l'état unique}

La structure discrète du continu ne s'obtient pas en assemblant cinq théories
indépendantes. Ses cinq aspects internes ---spectral, solénoïdal, de jauge,
cohomologique et déterminantal--- émergent corrélativement du même état
APP--TRIT--TPK. Cet ouvrage ne démontre pas de nouveau le traité intégral, mais
préserve le fait causal dont le radion a besoin : la section coinductive qui
génère \(\pi,e,\varphi\) et \(\alpha\) ne vit pas sur une droite réelle nue ; elle vit
dans un objet doté de fibres, de relations, d'orientation, de retenue, de mémoire, de terminal et de
lecteur.

\begin{statusbox}
\textbf{Falsificateur de généalogie.} Si une formule ultérieure reçoit
\(180/\pi\), \(\epsR\), une masse ou une table métrologique pour choisir une phase,
un coefficient ou une branche, elle cesse d'être une génération HMT-first. La confrontation
externe peut reconnaître une sortie ; elle ne peut pas la sélectionner.
\end{statusbox}


\subsection{Indépendance prospective de la clôture par rapport à la valeur finale}
% Fragmento público generado desde propietarios íntegros.
% Residencia: 52.1.2.
% Fuente: ../incoming/radion_monografia_20260827/manuscrito/sections/02_genealogia.tex:181-272

\subsubsection{Généalogie des constantes et construction du couple angulaire}

\paragraph{Section coinductive et unicité de la valeur générée de \(\pi\)}

Soit \(x_\infty\) la section globale du système projectif. Les lecteurs
coinductifs produisent
\[
\operatorname{ev}_\pi(x_\infty)=\pih,\qquad
\operatorname{ev}_e(x_\infty)=e_{\HMT},\qquad
\operatorname{ev}_\varphi(x_\infty)=\varphi_{\HMT}.
\]
L'indice HMT ne désigne pas un autre nombre. Il conserve la généalogie de la même
constante. En particulier,
\[
\pih=3.14159265358979323846264338327950288419716939937510\ldots.
\]
Le tour positif est
\begin{equation}
T_+(x_\infty)=2\pih.
\label{eq:vuelta-plus}
\end{equation}
Sa publication angulaire d'une unité est
\begin{equation}
R_{\HMT}^{\circ}=\frac{360^\circ}{T_+}=\frac{180^\circ}{\pih}.
\label{eq:radian-publicado}
\end{equation}
L'équation ne définit pas encore le radion enrichi ; elle définit la carte visible
du radian une fois le tour généré.

\paragraph{Bidirectionnalité des chemins et prolongations conjuguées}

La bidirectionnalité ne s'énonce pas comme un slogan. L'involution concrète
\begin{equation}
J_\alpha(T)=\frac{\alphah^2}{T},\qquad
T_-=J_\alpha(T_+),\qquad T_+T_-=\alphah^2
\label{eq:bidireccional-turn}
\end{equation}
conserve un produit et inverse le feuillet. Expansion et contraction sont deux
prolongations conjuguées du même état. La racine finie d'un jet
quadratique peut approcher \(T_+\), mais ne remplace pas la section
coinductive : la différence certifiée est
\[
T_{P_9}-T_+=-5.1110566290335\ldots\times10^{-25}.
\]
Le contrôle fini est un témoin ; le fondement est la limite.

\paragraph{Dépendance causale entre \(\alpha\), action et publication angulaire}

La clôture dodécaphasique génère \(\alphah\) à partir de l'état signé et de ses
cartes réversibles. Ensuite, le lecteur déterminantal local fixe la décennie
d'action au moyen de
\[
\Sigma_H=\varphi_{\HMT}\alphah^{16},
\qquad \operatorname{ord}_{10}\Sigma_H=-34.
\]
La section de retour produit \(\hret\). Ni \(\alpha\) ni \(\hbar\) ne se
\enquote{mesurent en degrés}. Ils sont publiés ultérieurement au moyen du couple angulaire
\begin{equation}
P(\alpha,\eta_{\mathrm{ret}})=
\left([10^3\alpha]_{360},[2(\eta_{\mathrm{ret}}+\alpha)]_{360}\right)
=(A,C^*).
\label{eq:parangular}
\end{equation}
Dans la branche canonique,
\[
A=7.297352569283800997285105472380663\ldots{}^\circ,
\]
\[
C^*=2.1237383389686457242619513803933607937\ldots{}^\circ.
\]
Alors seulement, la carte analytique
\[
x=\frac{\pi A}{180},\qquad y=\frac{\pi C^*}{180}
\]
ouvre les canaux \(q_\pm=e^{-x\mp y}\).

L’ordre causal complet est
\[
\alpha\to(\eta\parallel-34)\to\hbar
\to P\to\kappa_{360}\to q_\pm
\to\text{réalisations MD}.
\]
La relation \(C^*/2-\alpha\) peut servir de vérification inverse. Elle n'est pas utilisée
pour générer \(\hret\) rétroactivement.

\begin{narrativebox}
La carte angulaire ne transforme pas une constante adimensionnelle ou une action en un
angle par décret. Elle publie deux coordonnées de l'état sur une roue finie.
Le degré est le langage de cette publication ; il n'est pas la substance ontologique de
\(\alpha\) ni de \(\hbar\).
\end{narrativebox}
% Fuente: ../incoming/radion_monografia_20260827/manuscrito/sections/03_cierre_exacto.tex:254-285
\paragraph{Indépendance prospective par rapport à la valeur finale}

Le vérificateur de l'opérateur de profondeur interdit expressément trois entrées :
\[
\frac{180}{\pi},\qquad \epsR,\qquad
\text{le décimal cible}.
\]
Il construit d'abord \(S_E\) et \(S_T\), puis exécute la soustraction APP, et ne
compare la clôture qu'à la fin. Les contrôles donnent
\[
\texttt{target\_epsilon\_used=False},\qquad
\texttt{target\_180\_over\_pi\_used=False}.
\]
De même :
\begin{itemize}
\item un seul feuillet donne un défaut d'ordre \(10^{-5}\) ;
\item remplacer \(\theta_2\) par \(\theta_1\) ou \(\theta_3\) déplace le
résultat d'environ \(\pm\pi/6\) ;
\item éliminer le canal \(e\) laisse un défaut d'ordre \(10^{-3}\) ;
\item la variante historique \(\Delta_4(A,C^*)\) produit un autre nombre ;
\item la queue spectrale postérieure à \(\Rstate\) et l'opérateur harmonique à
quatre termes ne coïncident pas non plus avec \(\epsone\).
\end{itemize}

\begin{statusbox}
\textbf{Statut.} La clôture est une \status{formalisation nouvelle} démontrée
avec une structure de départ explicite et un certificat informatique indépendant de la valeur finale.
L'égalité est exacte à la limite coinductive. Le profil court de 36 chiffres
de \(\alpha\) génère une variante qui ne diverge qu'après de nombreux chiffres ;
il est conservé comme témoin historique, non comme valeur canonique.
\end{statusbox}



\section{Sections directe et torsionnelle : retenue APP et clôture exacte du radion}
\subsection{Construction indépendante des sections \(S_E\) et \(S_T\) avant la définition du cocycle \(\varepsilon_R=S_T-S_E\)}
% Fragmento público generado desde propietarios íntegros.
% Residencia: 52.2.1.
% Fuente: ../incoming/radion_monografia_20260827/manuscrito/sections/03_cierre_exacto.tex:1-50
\noindent La clôture exacte s'obtient au moyen de deux sections construites avant le cocycle qui les relie.\par\medskip

\subsubsection{Construction indépendante des sections directe et torsionnelle}

\paragraph{Section directe du système temporel dodécaphasique}

La roue dodécaphasique a pour centres
\begin{equation}
\theta_m=30m-10^\circ,\qquad m=1,\ldots,12.
\label{eq:rueda-dodecafasica}
\end{equation}
La deuxième phase est \(\theta_2=50^\circ\). Avec la publication
\(A=10^3\alphah\), elle forme la section directe
\begin{align}
S_E
&=\frac{T_+}{360}\,(50+10^3\alphah)\\
&=2\pih\left(\frac5{36}+\frac{25}{9}\alphah\right).
\label{eq:seccion-electrica}
\end{align}
L'indice \(E\) signifie ici \emph{directe ou commune} ; sa réalisation
électrique viendra après le lecteur constitutif. On ne postule pas
littéralement \(A=E\).

Les cylindres coinductifs prouvent
\begin{equation}
1<S_E<2.
\label{eq:SE-intervalo}
\end{equation}
La division APP impose alors
\begin{equation}
q_E=\lfloor S_E\rfloor=1,
\qquad
D_E=\operatorname{rem}_1(S_E)=S_E-1.
\label{eq:DE}
\end{equation}
L'entier un n'est pas le nombre de tours de l'hélice. C'est le quotient unique
de la section directe au sein de la cellule d'unité.

\paragraph{Section torsionnelle déterminée par l'opérateur \(\Delta_4\)}

La seconde section utilise une sortie coinductive antérieure :
\begin{equation}
\Delta_4=
\frac{\pih-e_{\HMT}\log\pih}{270}
\left(1+\frac{\pih}{729}\right).
\label{eq:delta4}
\end{equation}
Le mot \emph{torsion} désigne ici le quantum global de la carte HMT. Il ne
l'identifie pas automatiquement à la torsion de Cartan
\(T^I=D_\omega e^I\). Le pont vers Cartan exige une autre flèche typée.

\subsection{Densité orientée du secteur actif et détermination unique de \(20/81\)}
% Fragmento público generado desde propietarios íntegros.
% Residencia: 52.2.2.
% Fuente: ../incoming/radion_monografia_20260827/manuscrito/sections/03_cierre_exacto.tex:51-103

\paragraph{Densité orientée \(20/81\) du secteur actif \(90/120\)}

Le sélecteur tritique de phase sépare
\[
H_+=\{1,4,7\},\qquad H_-=\{2,5,8\},\qquad H_0=\{3,6,9\}.
\]
Il y a six phases actives et quinze paires non ordonnées :
\[
H_{\mathrm{act}}=H_+\cup H_-,\qquad
\left|\binom{H_{\mathrm{act}}}2\right|=\binom62=15.
\]
Le mode actif et son extension à huit positions sont
\begin{equation}
N_6=6\binom62=90,\qquad
N_8=8\binom62=120.
\label{eq:90-120-incidencia}
\end{equation}
Deux feuillets conjugués du secteur actif, normalisés par la fibre ambiante,
donnent
\begin{equation}
\rho_{\mathrm{tw}}
=\frac{2N_6}{|\F_3^6|}
=\frac{2\cdot90}{729}
=\frac{180}{729}
=\boxed{\frac{20}{81}}.
\label{eq:rho-twoway}
\end{equation}

\begin{proposition}[Unicité relative de la densité]
Le secteur actif \(N_6=90\), les deux feuillets conjugués et la
normalisation par \(|\F_3^6|=729\) étant fixés, la densité orientée est nécessairement
\(20/81\).
\end{proposition}
\begin{proof}
La cardinalité transportée est \(2N_6=180\). La normalisation par l'espace ambiant
donne \(180/729\) ; diviser le numérateur et le dénominateur par \(9\) produit \(20/81\).
Aucune comparaison avec \(180/\pi\) ni avec \(\epsR\) n'intervient.
\end{proof}

La lecture historique comme projection sur un réseau \(1/81\) est un contrôle de
rigidité utile, mais n'est plus le fondement : le coefficient s'obtient avant
la clôture. L'ablation d'un feuillet produit \(10/81\), et la sortie change à
l'ordre \(10^{-5}\) ; le facteur deux n'est donc pas décoratif.

La section torsionnelle est
\begin{equation}
\Phi_T=\rho_{\mathrm{tw}}\Delta_4,
\qquad
S_T=1+\Phi_T.
\label{eq:seccion-torsional}
\end{equation}


\subsection{La retenue \(\varepsilon_R\) comme cocycle de transition entre sections}
% Fragmento público generado desde propietarios íntegros.
% Residencia: 52.2.3.
% Fuente: ../incoming/radion_monografia_20260827/manuscrito/sections/03_cierre_exacto.tex:104-253
\subsubsection{Cocycle de transition entre les sections directe et torsionnelle}

\paragraph{Définition opératoire de la retenue \(\varepsilon_R\)}

\begin{definition}[Retenue réelle du radion]
La sortie unitaire du radion est le cocycle de transition
\begin{equation}
\epsone:=S_T-S_E
=\Phi_T-D_E
=\frac{20}{81}\Delta_4-operatorname{rem}_1(S_E).
\label{eq:epsilon-def}
\end{equation}
\end{definition}

L'équation contient une soustraction, mais ce n'est pas un ajustement à la cible. Ses deux
opérandes \(S_T\) et \(S_E\) sont générés séparément. L'indépendance
démontrée est une indépendance par rapport à \(180/\pi\) et à un
\(\varepsilon\) cible ; on n'affirme pas d'indépendance algébrique entre une
différence et ses termes.

\begin{theorem}[Clôture exacte unitaire]
Avec \(S_E\), \(\Phi_T\) et \(\epsone\) définis par
\eqref{eq:seccion-electrica}, \eqref{eq:delta4}, \eqref{eq:rho-twoway} et
\eqref{eq:epsilon-def}, on a exactement
\begin{equation}
\boxed{1=S_E-\Phi_T+\epsone.}
\label{eq:cierre-unitario}
\end{equation}
\end{theorem}
\begin{proof}
Par \eqref{eq:SE-intervalo}, \(D_E=S_E-1\). Par définition,
\(\epsone=\Phi_T-D_E\). Alors
\[
S_E-\Phi_T+\epsone
=S_E-\Phi_T+\Phi_T-(S_E-1)=1.
\]
La force de la preuve ne réside pas dans cette annulation élémentaire, mais dans le fait que la
généalogie préalable fixe les deux opérandes sans utiliser la clôture comme entrée.
\end{proof}

\paragraph{Publication de la clôture dans la carte angulaire}

La carte du tour multiplie une fraction unitaire par
\(360/T_+=180/\pih\). En particulier,
\[
\epsR^\circ=\frac{360^\circ}{T_+}\epsone.
\]
Multiplier \eqref{eq:cierre-unitario} par \(360^\circ/T_+\) et utiliser
\[
\frac{360}{T_+}S_E
=360\left(\frac5{36}+\frac{25}{9}\alphah\right)
=50+1000\alphah=50+A
\]
produit la clôture angulaire.

\begin{theorem}[Équation exacte du radion]
\BigRadionEquation
\end{theorem}

L'équation dit quelque chose de plus précis que \enquote{le radian est
\(57.2957\ldots{}^\circ\)}. Elle décompose cette publication en quatre opérations
avec démonstration de référence :
\begin{center}
\begin{tabularx}{.96\textwidth}{>{\bfseries}l X}
\toprule
\(50^\circ\) & deuxième phase de la roue ; charnière critique dans une carte déclarée ;\\
\(A^\circ\) & coordonnée commune de la publication parangulaire ;\\
\(-\frac{20}{81}\Delta_4^\circ\) & transport orienté de la torsion globale ;\\
\(+\epsR^\circ\) & retenue réelle qui raccorde les deux relèvements.\\
\bottomrule
\end{tabularx}
\end{center}

\paragraph{Convergence selon la profondeur du raffinement}

À la profondeur \(N\), le générateur fournit des cylindres rationnels
\(p_N\to\pih\) et \(e_N\to e_{\HMT}\). On définit
\[
S_{E,N}=2p_N\left(\frac5{36}+\frac{25}{9}\alphah\right),
\]
\[
\Delta_{4,N}=\frac{p_N-e_N\log p_N}{270}
\left(1+\frac{p_N}{729}\right),\qquad
S_{T,N}=1+\frac{20}{81}\Delta_{4,N},
\]
et
\[
\mathcal E_N=S_{T,N}-S_{E,N}.
\]
La continuité de la fonctionnelle en \(3<p<4\), \(2<e<3\), jointe aux
dérivées monotones, transporte chaque couple de cylindres vers un intervalle
certifié. À la profondeur de \(45\) blocs, sa largeur est inférieure à
\(4.81\times10^{-130}\).

Chaque retour de neuf niveaux contient six trits par niveau. C'est pourquoi le taux
de raffinement est
\begin{equation}
729^{-9}=3^{-54}.
\label{eq:contraccion-54}
\end{equation}
Ce chiffre contrôle la largeur des intervalles ; ce n'est pas la valeur de
\(\epsR\).

La soustraction de triades stabilise
\begin{center}
\ttfamily
000|000|000|896|802|822|044|562|981|999|050|695|020|651|286|413
\end{center}
et le préfixe de retenues
\begin{center}
\ttfamily
0,0,0,0,0,-1,0,-1,-1,-1,0,-1,0.
\end{center}
Tel est le mécanisme concret de mémoire décimale : non une queue choisie, mais
la récurrence \eqref{eq:subacarreo} appliquée à des opérandes préalablement
construits.

\paragraph{Évaluation des valeurs canoniques}

La sortie coinductive est
\begin{align}
\epsone
&=8.9680282204456298199905069502065128641372736205009\ldots
\times10^{-10},\\
\epsR^\circ
&=5.1383016758575282071685164622645947489908574400054\ldots
\times10^{-8}\,{}^\circ.
\label{eq:epsilon-valores}
\end{align}
La publication complète donne
\[
\frac{180^\circ}{\pih}
=57.2957795130823208767981548141051703324054724665643\ldots{}^\circ.
\]

\begin{exactbox}
\begin{align*}
50^\circ+A^\circ
&=57.29735256928380099728510547238066299956\ldots{}^\circ,\\
\frac{20}{81}\Delta_4^\circ
&=0.00157310758449687906223272996065728980465\ldots{}^\circ,\\
50^\circ+A^\circ-\frac{20}{81}\Delta_4^\circ
&=57.29577946169930411822287274242000570976\ldots{}^\circ,\\
\epsR^\circ
&=0.00000005138301675857528207168516462264594749\ldots{}^\circ,\\
\text{somme finale}
&=57.29577951308232087679815481410517033241\ldots{}^\circ.
\end{align*}
\end{exactbox}



\section{Coordonnée critique et système temporel dodécaphasique orienté : publication affine de la phase de \(50^\circ\)}
\subsection{Application affine \(j_{100}\) de la carte APP et unicité de la phase \(m=2\) dans \(j_{100}(1/2)=\theta_m=50^\circ\)}
% Fragmento público generado desde propietarios íntegros.
% Residencia: 52.3.1.
% Fuente: ../incoming/radion_monografia_20260827/manuscrito/sections/04_cincuenta_critica.tex:1-105
\noindent La coordonnée critique relie la carte de Cayley, le registre APP et la phase dodécaphasique par des applications typées.\par\medskip

\subsubsection{Deux réalisations exactes de la coordonnée angulaire \(50^\circ\)}

\paragraph{Deuxième phase du système temporel dodécaphasique}

Par \eqref{eq:rueda-dodecafasica}, la suite des centres commence
\[
20^\circ,\ 50^\circ,\ 80^\circ,\ 110^\circ,\ldots.
\]
Résoudre
\[
30m-10=50
\]
donne uniquement \(m=2\). Par conséquent, \(50^\circ\) n'est ni un arrondi de \(A\),
ni un demi-tour ni le centre d'un cercle : c'est une phase typée du calendrier
dodécaphasique. La fraction de tour est
\[
\frac{50}{360}=\frac5{36},
\]
exactement le premier terme de \(S_E/T_+\).

\paragraph{Égalisation de la coordonnée critique de Cayley}

L'involution fonctionnelle
\[
\rho\longmapsto1-\bar\rho
\]
fixe
\[
\operatorname{Fix}(\rho\mapsto1-\bar\rho)
=\{\rho:\Re\rho=1/2\}.
\]
Le nombre \(1/2\) est la coordonnée réelle autoduale de l'intervalle fonctionnel
\(0\leq\Re s\leq1\). Ce n'est pas \enquote{la moitié d'une droite} au sens
longitudinal, et il ne contient pas la hauteur spectrale \(\gamma\). Pour
\[
\rho_\gamma=\frac12+i\gamma,
\]
la transformation de Cayley peut s'écrire
\begin{equation}
\tau_\gamma
=\frac{\rho_\gamma-1}{\rho_\gamma}
=\frac{\gamma+i/2}{\gamma-i/2}
\in S^1,
\qquad
\gamma=\frac12\cot\frac{\theta_\gamma}{2}.
\label{eq:cayley-critica}
\end{equation}
Ainsi, \(1/2\) fixe la couture, \(\gamma\) détermine le caractère angulaire et
\(\Gamma_9\) conserve la profondeur de ses itérations.

\paragraph{Injection affine \(j_{100}\) du registre APP dans la carte angulaire}

La complétion cubique dispose une chambre de dix marques par coordonnée. Une
face contient \(10^2=100\) positions. Nous déclarons la carte affine
\begin{equation}
j_{100}:[0,1]\longrightarrow[0,100^\circ],
\qquad j_{100}(\sigma)=100\sigma^\circ.
\label{eq:j100}
\end{equation}
Alors
\[
j_{100}\left(\frac12\right)=50^\circ.
\]
En exigeant simultanément \(j_{100}(1/2)=\theta_m\), il vient
\[
100\cdot\frac12=30m-10
\quad\Longleftrightarrow\quad m=2.
\]

\begin{theorem}[Égalisateur critique--dodécaphasique]
Dans la carte APP de cent marques, la coordonnée fixe de l'involution critique
se publie exactement dans la deuxième phase de la roue :
\begin{equation}
\boxed{j_{100}(1/2)=\theta_2=50^\circ.}
\label{eq:igualador-50}
\end{equation}
La phase \(m=2\) est unique.
\end{theorem}

Le théorème est exact une fois \eqref{eq:j100} déclarée. Son statut est
\status{formalisation nouvelle}: le demi critique, la complétion \(10^3\), la
face de cent marques et \(\theta_2\) étaient déjà construits ; le choix de cette
face comme carte critique affine réunit pour la première fois leurs domaines.

\paragraph{Compatibilité de la valeur critique avec la carte de Cayley}

La transformation de Cayley n'envoie pas \(1/2\) sur \(50^\circ\). Si
\(\gamma=0\), \eqref{eq:cayley-critica} donne
\[
\tau_0=-1=e^{i\pi},
\]
c'est-à-dire \(180^\circ\). Les \(50^\circ\) appartiennent à la carte APP de cent
marques et au calendrier dodécaphasique ; l'angle de Cayley dépend de \(\gamma\).
Cette distinction empêche de réduire tous les zéros de la fonction zêta à une seule
phase.

\begin{falsifier}
L'identification \(1/2\leftrightarrow50^\circ\) échoue si l'on élimine la
carte \eqref{eq:j100}. L'égalité scalaire \(100(1/2)=50\) ne suffit pas
à identifier à elle seule deux objets. Le théorème dépend de la face APP
distinguée et de la roue dodécaphasique.
\end{falsifier}


\subsection{Compatibilité de la phase critique avec le retour nonadique et l'avance de mémoire}
% Fragmento público generado desde propietarios íntegros.
% Residencia: 52.3.2.
% Fuente: ../incoming/radion_monografia_20260827/manuscrito/sections/04_cincuenta_critica.tex:106-171
\subsubsection{Séparation typée des 2 objets de valeur \(1/2\)}

APP contient en outre un mode singulier avec
\[
s=\frac12,\qquad s^2=\frac14,\qquad
1-s^2=\frac34=\frac{90}{120},\qquad
s\sqrt{1-s^2}=\frac{\sqrt3}{4}.
\]
Ce \(s\) naît des angles principaux des feuillets somme et produit. La
couture critique \(\Re\rho=1/2\) naît d'une involution fonctionnelle. Que les deux
coordonnées valent \(1/2\) est une coïncidence scalaire dotée d'un contenu possible,
mais n'autorise pas à les identifier sans un carré typé.

La manière correcte de conserver la découverte est un diagramme d'égalisateurs :
\[
\begin{tikzpicture}[baseline=(current bounding box.center),node distance=14mm and 23mm,
 every node/.style={font=\small}]
\node[draw=RadBlue,rounded corners,fill=RadMist] (crit) {couture critique \(\Re s=1/2\)};
\node[draw=RadTeal,rounded corners,fill=RadCream,right=of crit] (face) {carte \(j_{100}\)};
\node[draw=RadCopper,rounded corners,fill=RadCream,right=of face] (theta) {phase \(\theta_2=50^\circ\)};
\node[draw=RadRose,rounded corners,fill=white,below=of face] (mode) {mode APP \(s=1/2\), \(1-s^2=90/120\)};
\draw[-{Latex},thick,RadBlue] (crit)--node[above]{\(j_{100}\)}(face);
\draw[-{Latex},thick,RadTeal] (face)--node[above]{égalisateur}(theta);
\draw[-{Latex},dashed,RadRose] (mode)--node[right,align=left]{même chiffre;\\domaine distinct}(face);
\end{tikzpicture}
\]
Le diagramme ne simule pas la flèche absente. Il montre ce qui est démontré, ce qui a
été formalisé et quelle égalité nécessite encore un transport supplémentaire si
l'on vise une identification opératoire globale.

\subsubsection{Compatibilité entre cercle spectral et mémoire hélicoïdale}

Le lecteur de Cayley--Pontryagin transporte une fibre spectrale au moyen de
\[
\widetilde U_{\Gamma_9^n}J\psi_\gamma
=\tau_\gamma^nJ\psi_\gamma.
\]
La publication circulaire est \(\tau_\gamma^n\in S^1\). Le relèvement
\[
(\tau_\gamma^n,n)\in S^1\times\Z
\]
conserve le nombre de retours. Sur le feuillet contractant apparaît
\[
\widetilde M_9^nJ\psi_\gamma
=\left(\frac59\right)^n\tau_\gamma^nJ\psi_\gamma.
\]
Nous avons donc trois images typées :
\[
\text{cercle de phase},\qquad
\text{hélice de mémoire},\qquad
\text{spirale contractante}.
\]

Un zéro n'est pas un tour, et un neuf n'est pas un zéro particulier de la fonction
zêta. Le zéro fixe un caractère \(\tau_\gamma\) qui persiste à travers les
tours ; le retour nonadique incrémente la mémoire. Le neuf fonctionne comme zéro
résiduel dans \(\Z/9\Z\), opération arithmétique distincte de l'énumération
des zéros spectraux.

\begin{narrativebox}
La droite critique devient un cercle lorsque nous oublions la hauteur linéaire et
conservons le caractère unitaire. Elle devient une hélice lorsque nous réinsérons
la mémoire du nombre de fois où le caractère a été transporté. La profondeur
n'est pas une troisième coordonnée décorative : c'est l'information que la projection
circulaire ne peut pas retenir.
\end{narrativebox}


\section{Couple angulaire et action réduite : construction de l'ellipse positive et de l'aire par radion}
\subsection{Loi surfacique du radion : \(\mathcal A(\theta)=\hbar_{\mathrm{ret}}\theta\) et \(\mathcal A(2\pi)=h_{\mathrm{ret}}\)}
% Fragmento público generado desde propietarios íntegros.
% Residencia: 52.4.1.
% Fuente: ../incoming/radion_monografia_20260827/manuscrito/sections/05_geometria_accion.tex:1-217
\noindent La géométrie spectrale détermine la forme et la loi surfacique détermine l'action transportée par la phase.\par\medskip

\subsubsection{Équivalence entre les cartes spectrale et géométrique}

\paragraph{Opérateur spectral positif et valeurs propres conjuguées}

Le couple angulaire produit l'opérateur bicouche
\begin{equation}
B_1=AI+C^*R,\qquad R^2=I,\qquad
P_\pm=\frac{I\pm R}{2}.
\label{eq:B1}
\end{equation}
Ses valeurs propres sont \(A\pm C^*\). Comme \(0<C^*<A\), les deux sont positives.
Dans la carte \(\kappa=\pi/180\), l'ellipse spectrale a pour demi-axes
\begin{equation}
a_1^2=\kappa(A+C^*),\qquad
b_1^2=\kappa(A-C^*).
\label{eq:semiejes-espectrales}
\end{equation}
La décomposition est réversible :
\begin{equation}
A=\frac{a_1^2+b_1^2}{2\kappa},\qquad
C^*=\frac{a_1^2-b_1^2}{2\kappa}.
\label{eq:reconstruccion-AC}
\end{equation}
L'ellipse n'est pas dessinée pour illustrer un nombre. Elle reconstruit exactement le
couple qui l'a générée.

Définissons
\begin{equation}
\rho_1=\sqrt{A^2-C^{*2}},\qquad
\eta=\operatorname{artanh}\frac{C^*}{A}
=\frac12\log\frac{A+C^*}{A-C^*}.
\label{eq:eta}
\end{equation}
Alors
\[
A\pm C^*=\rho_1e^{\pm\eta},
\]
et
\begin{equation}
a_1=\sqrt{\kappa\rho_1}\,e^{\eta/2},
\qquad
b_1=\sqrt{\kappa\rho_1}\,e^{-\eta/2}.
\label{eq:semiejes-eta}
\end{equation}
Le produit conserve l'échelle ; le quotient conserve la forme :
\[
a_1b_1=\kappa\rho_1,
\qquad
\frac{b_1}{a_1}=e^{-\eta}
=\sqrt{\frac{A-C^*}{A+C^*}}.
\]

\paragraph{Invariants focaux de l'opérateur positif}

La demi-distance focale est
\begin{equation}
c_1^2=a_1^2-b_1^2=2\kappa C^*,
\qquad
F_\pm=(\pm c_1,0).
\label{eq:focos}
\end{equation}
La distance totale vaut
\begin{equation}
d_1=2c_1=2\sqrt{2\kappa C^*}
=0.544545509325626623406299779866023\ldots.
\label{eq:d1}
\end{equation}
Le nombre isolé n'est pas la thèse. L'identité \(c_1^2=2\kappa C^*\) dit que
\(C^*\) est littéralement la scission focale quadratique dans la carte
distinguée. Les foyers sont les extrémités conjuguées de l'ouverture spectrale ;
ils ne sont pas, à eux seuls, les pôles électrique et magnétique. Cette lecture nécessite
le transducteur constitutif.

L'excentricité est
\begin{equation}
e_f=\frac{c_1}{a_1},\qquad
e_f^2=1-e^{-2\eta}=\frac{2C^*}{A+C^*}.
\label{eq:excentricidad}
\end{equation}
Numériquement,
\[
\begin{aligned}
\eta&=0.299689683921630799684926562863927\ldots,\\
e_f&=0.671451895550191986916277055864332\ldots.
\end{aligned}
\]

\paragraph{Bloc APP de référence et transport normalisé}

Le bloc central
\[
B_0=7I+2R=9P_++5P_-
\]
possède
\[
\rho_0=\sqrt{45},\qquad
\eta_0=\frac12\log\frac95,\qquad
e_0=\frac23.
\]
Sa distance focale est
\[
d_0=4\sqrt\kappa
=0.528443639680801468446003835349358\ldots.
\]
La comparaison exacte est
\begin{equation}
\boxed{\left(\frac{d_1}{d_0}\right)^2=\frac{C^*}{2}},
\qquad
d_1=d_0\sqrt{\frac{C^*}{2}}.
\label{eq:invariante-focal-relativo}
\end{equation}
Telle est la généalogie défendable de \(0.544545\ldots\) : une échelle focale APP
de référence multipliée par un invariant relatif. Sa proximité visuelle avec
\(0.54\) n'est pas une identité avec le défaut APP \(54\).

Le transport depuis \(B_0\) se sépare en échelle et anisotropie :
\begin{equation}
T_{\mathrm{rel}}=\frac{\rho_1}{\sqrt{45}}
\exp\bigl((\eta-\eta_0)R\bigr).
\label{eq:transporte-rel}
\end{equation}
Sur les feuillets,
\[
t_+=\frac{A+C^*}{9}=1.0467878786947163\ldots,
\qquad
t_-=\frac{A-C^*}{5}=1.0347228460630311\ldots.
\]
La déformation n'est pas une seule décimale. Sa partie isotrope est
\(\sqrt{t_+t_-}\) ; sa partie hyperbolique est
\(\tfrac12\log(t_+/t_-)\).

\subsubsection{Ellipse positive associée au couple angulaire}

\paragraph{Construction au moyen de demi-axes conjugués}

L'échelle spectrale précédente est dépourvue d'unités d'action. La mécanique
dimensionnelle publie la même forme sur la droite interne d'action :
\begin{equation}
a_S=\sqrt{2\hret e^\eta},
\qquad
b_S=\sqrt{2\hret e^{-\eta}}.
\label{eq:semiejes-accion}
\end{equation}
Immédiatement,
\begin{equation}
a_Sb_S=2\hret,\qquad
\frac{a_S}{b_S}=e^\eta.
\label{eq:producto-forma}
\end{equation}
Le produit fixe l'action ; le quotient fixe l'anisotropie. L'ellipse est
paramétrée par
\begin{equation}
Q(\theta)=a_S\cos\theta,\qquad
P(\theta)=b_S\sin\theta.
\label{eq:elipse-param}
\end{equation}

Les ellipses spectrale et d'action sont homothétiques :
\begin{equation}
\lambda_{\mathrm{act}}^2=\frac{2\hret}{\kappa\rho_1},
\qquad
(a_S,b_S,c_S)=\lambda_{\mathrm{act}}(a_1,b_1,c_1).
\label{eq:homotecia}
\end{equation}
Elles partagent la forme \(e^{-\eta}\), non la dimension physique.

\paragraph{Loi surfacique de l'action par radion}

La forme symplectique canonique est \(\omega=\diff Q\wedge\diff P\). L'aire
orientée balayée de \(0\) à \(\theta\) est
\begin{align}
\mathcal A(\theta)
&=\frac12\int_0^\theta
\bigl(Q(t)P'(t)-P(t)Q'(t)\bigr)\,\diff t\\
&=\frac12a_Sb_S\int_0^\theta
(\cos^2t+\sin^2t)\,\diff t.
\end{align}
En utilisant \eqref{eq:producto-forma}, nous obtenons le résultat central.

\begin{theorem}[Action par radion]
Pour l'ellipse \eqref{eq:elipse-param},
\begin{equation}
\boxed{\mathcal A(\theta)=\hret\theta.}
\label{eq:area-theorem}
\end{equation}
En particulier,
\begin{equation}
\boxed{\mathcal A(1)=\hret},
\qquad
\boxed{\mathcal A(2\pi)=2\pi\hret=\hfull}.
\label{eq:area-radion-vuelta}
\end{equation}
\end{theorem}

Cette égalité permet de définir le radion sans faire d'abord appel à la longueur
d'arc.

\begin{definition}[Radion orienté]
\begin{equation}
\mathfrak r_\sigma=(\theta=1,\sigma),
\qquad \sigma\in\{+1,-1\},
\qquad
\mathcal A(\mathfrak r_\sigma)=\sigma\hret.
\label{eq:radion-oriented}
\end{equation}
\end{definition}

Le radian est la projection qui oublie \(\sigma\), le feuillet, la mémoire et
l'aire. Le radion conserve cette information. D'où le nom \emph{rad--ion} :
l'orientation précède la charge physique, et la charge apparaît plus tard comme
\[
Q(x)=n_Q(x)e_{\mathrm{int}},\qquad
e_{\mathrm{int}}=\sqrt{\frac{2\alphah\hfull}{Z_{\mathrm{int}}}}.
\]


\subsection{Demi-axes conjugués, invariants focaux et orientation de l'anisotropie}
% Fragmento público generado desde propietarios íntegros.
% Residencia: 52.4.2.
% Fuente: ../incoming/radion_monografia_20260827/manuscrito/sections/05_geometria_accion.tex:218-290
\paragraph{Correspondance entre phase, action, aire et orientation}

\begin{figure}[H]
\centering
\begin{tikzpicture}[scale=1.12]
  \def\aa{4.0}
  \def\bb{2.95}
  \def\cc{2.70}
  \draw[->,RadGray] (-4.7,0)--(4.7,0) node[right]{\(Q\)};
  \draw[->,RadGray] (0,-3.5)--(0,3.5) node[above]{\(P\)};
  \fill[RadTeal!8] (0,0) ellipse[x radius=\aa,y radius=\bb];
  \draw[RadTeal,line width=1.35pt] (0,0) ellipse[x radius=\aa,y radius=\bb];
  \fill[RadCopper] (-\cc,0) circle (2pt) node[below=2mm]{\(F_-\)};
  \fill[RadCopper] (\cc,0) circle (2pt) node[below=2mm]{\(F_+\)};
  \draw[RadGold,line width=1.1pt,-{Latex}] (\aa,0)
     arc[start angle=0,end angle=57.3,x radius=\aa,y radius=\bb];
  \node[RadGold] at (3.7,2.3) {un radion};
  \draw[dashed,RadBlue] (0,0)--(\aa,0) node[midway,below]{\(a_S\)};
  \draw[dashed,RadBlue] (0,0)--(0,\bb) node[midway,left]{\(b_S\)};
  \node[align=center,RadNavy] at (0,-3.25)
    {aire totale \(\pi a_Sb_S=2\pi\hret=\hfull\)};
\end{tikzpicture}
\caption{L'ellipse d'action. Un arc paramétrique d'un radian balaie
\(\hret\) ; un tour enferme \(\hfull\). Les foyers codent l'ouverture
\(C^*\) dans la carte distinguée.}
\label{fig:elipse-accion}
\end{figure}

\paragraph{Relation entre l'action réduite, l'incertitude et l'aire}

La relation \(\hbar=h/(2\pi)\) est souvent présentée comme une commodité
de notation : lorsqu'on décrit des oscillations, des rotations et des transformées de Fourier,
les facteurs \(2\pi\) disparaissent si l'on utilise \(\hbar\). L'ellipse d'action
révèle une lecture plus forte. La division par \(2\pi\) est exactement le
passage de l'action d'un tour à l'action par unité de phase :
\[
h=\mathcal A(2\pi),\qquad \hbar=\mathcal A(1).
\]
La constante réduite n'est pas seulement \(h\) avec une notation plus commode ; c'est la
densité surfacique d'action par rapport à la phase.

Heisenberg, Dirac et Jordan ont établi le langage des variables conjuguées,
des commutateurs et des transformations canoniques. On ne leur attribue pas ici une
\enquote{ellipse ontologique} littérale sans source. La contribution HMT--MD consiste à
identifier la figure positive, à construire ses demi-axes à partir de \((A,C^*)\) et à
démontrer \eqref{eq:area-theorem}. La généalogie historique ouvre le langage ;
le théorème actuel fixe l'objet.

\begin{narrativebox}
L'action ne s'ajoute pas ensuite au mouvement. Dans l'ellipse du radion,
l'action est l'aire que le mouvement génère. Un tour ne transporte pas une
constante préfabriquée : il l'enferme.
\end{narrativebox}

\Needspace{10\baselineskip}
\paragraph{Contrôles algébriques des invariants focaux}

Les ressemblances décimales sont soumises à un contrôle :
\begin{center}
\begin{tabularx}{.96\textwidth}{l r X}
\toprule
Comparaison & Résidu & Verdict\\
\midrule
\(d_1\) face à \(0.54\) & \(4.5455\times10^{-3}\) & pas d'égalité ; le \(54\) ne se déduit pas ;\\
\(d_1^2\) face à \(8/27\) & \(2.3352\times10^{-4}\) & proximité sans application ;\\
\(\eta\) frente a \(3/10\) & \(-3.1032\times10^{-4}\) & no igualdad;\\
\(e_f\) face à \(2/3\) & \(4.7852\times10^{-3}\) & \(2/3\) appartient exactement à \(B_0\) ;\\
\(e^{-\eta}\) face à \(20/27\) & \(3.0740\times10^{-4}\) & no igualdad.\\
\bottomrule
\end{tabularx}
\end{center}
Les invariants solides sont \eqref{eq:eta}, \eqref{eq:excentricidad} et
\eqref{eq:invariante-focal-relativo}. Le critère HMT ne consiste pas à récompenser
une ressemblance décimale, mais à exiger objet, opération et conservation.


\section{Retour nonadique et intérieur projectif : disque de Klein et profondeur hélicoïdale de mémoire}
\subsection{Normalisation au disque et intérieur hyperbolique de Hilbert--Beltrami--Klein}
% Fragmento público generado desde propietarios íntegros.
% Residencia: 52.5.1.
% Fuente: ../incoming/radion_monografia_20260827/manuscrito/sections/06_interior_helice.tex:1-108
\noindent La géométrie projective distingue l'action finie de la profondeur hyperbolique et de la mémoire du retour.\par\medskip

\subsubsection{Intérieur hyperbolique de Hilbert--Beltrami--Klein}

\paragraph{Normalisation projective dans le disque de Klein}

L'application affine
\[
u=\frac{Q}{a_S},\qquad v=\frac{P}{b_S}
\]
envoie l'ellipse d'action sur le disque unité
\[
D=\{(u,v):u^2+v^2<1\}.
\]
Sur chaque corde, la distance de Hilbert est définie par le birapport
de ses extrémités de frontière. La métrique infinitésimale résultante est la forme
de Beltrami--Klein
\begin{equation}
\diff s_K^2=
\frac{(1-r^2)(\diff u^2+\diff v^2)+(u\diff u+v\diff v)^2}
{(1-r^2)^2},
\qquad r^2=u^2+v^2.
\label{eq:klein-metric}
\end{equation}
Avec cette normalisation, la courbure gaussienne est constante :
\begin{equation}
K=-1.
\label{eq:klein-curvature}
\end{equation}
Les géodésiques sont des cordes euclidiennes, mais leur longueur hyperbolique diverge
à l'approche de la frontière.

\paragraph{Finitude de la frontière et infinitude de la profondeur hyperbolique}

Le déterminant métrique de \eqref{eq:klein-metric} produit
\begin{equation}
\diff A_K=\frac{\diff u\,\diff v}{(1-r^2)^{3/2}}.
\label{eq:klein-area-density}
\end{equation}
L'aire du sous-disque \(r<R<1\) est
\begin{align}
A_K(R)
&=\int_0^{2\pi}\int_0^R
\frac{r\,\diff r\,\diff\theta}{(1-r^2)^{3/2}}\\
&=2\pi\left(\frac1{\sqrt{1-R^2}}-1\right).
\label{eq:klein-area-R}
\end{align}
\Needspace{4\baselineskip}
Par conséquent,
\[
\lim_{R\to1^-}A_K(R)=+\infty.
\]
Simultanément, la même frontière enferme l'aire symplectique
\[
\pi a_Sb_S=2\pi\hret=\hfull<\infty.
\]

\begin{theorem}[Finitude de l'action et infinitude de la profondeur]
L'ellipse du radion possède une action de tour finie \(\hfull\), tandis que
son intérieur, muni de la métrique de Hilbert--Beltrami--Klein, a une aire
hyperbolique infinie :
\begin{equation}
\boxed{
\operatorname{Area}_{\mathrm{simp}}(\partial\mathbb E_S)=\hfull<\infty,
\qquad
\operatorname{Area}_{K}(\mathbb E_S)=+\infty.}
\label{eq:finite-infinite}
\end{equation}
\end{theorem}

Il n'y a pas de contradiction : \(\operatorname{Area}_{\mathrm{simp}}\) et
\(\operatorname{Area}_K\) sont des mesures différentes. La première compte l'action
orientée en phase ; la seconde mesure la profondeur projective intérieure. La thèse
\enquote{bord fini, intérieur infini} est désormais une égalité et une limite,
non une métaphore.

\paragraph{Coordonnées focales dans l'intérieur de Klein}

Dans le disque normalisé, les foyers se situent en \((\pm e_f,0)\). La distance
du centre à un foyer est
\begin{equation}
u_f=\operatorname{artanh}e_f
=\operatorname{arcosh}(e^\eta)
=0.813382331175342333760889731088228\ldots.
\label{eq:klein-focal}
\end{equation}
La distance foyer--foyer est
\[
d_K(F_-,F_+)=2u_f
=1.62676466235068466752177946217646\ldots.
\]
On a
\begin{equation}
e_f=\tanh u_f,\qquad
e^{-\eta}=\operatorname{sech}u_f,\qquad
\eta=\log\cosh u_f.
\label{eq:focal-hyperbolic-chain}
\end{equation}
La région jusqu'au rayon focal a pour aire
\[
A_K(e_f)=2\pi(e^\eta-1)
=2.19559620884061869541206115980360\ldots.
\]

Toutes les ellipses ouvertes sont projectivement isométriques en tant que domaines de
Hilbert. C'est pourquoi la courbure \(-1\) seule ne mémorise pas l'excentricité. La
forme HMT réside dans la structure conjointe : carte affine distinguée, forme
symplectique, étiquetage des feuillets et paramétrisation de frontière.

\subsection{Différence entre retour de phase et retour de l'état enrichi}
% Fragmento público generado desde propietarios íntegros.
% Residencia: 52.5.2.
% Fuente: ../incoming/radion_monografia_20260827/manuscrito/sections/06_interior_helice.tex:109-201
\subsubsection{Relèvement hélicoïdal du retour nonadique}

\paragraph{Décomposition nonadique et quotient de mémoire}

L'opération élémentaire
\begin{equation}
n=9q_9(n)+\rho_9(n),
\qquad
H_9(n)=(\rho_9(n),q_9(n))
\label{eq:H9}
\end{equation}
satisfait
\begin{equation}
H_9(n+9)=H_9(n)+(0,1).
\label{eq:H9-helix}
\end{equation}
La première coordonnée se ferme ; la seconde s'élève. Cette identité est la
version arithmétique minimale d'une vis.

Pour une fibre spectrale,
\[
n\longmapsto(\tau_\gamma^n,n)
\]
est une hélice sur \(S^1\). Ajouter la contraction \((5/9)^n\) produit une
spirale intérieure. Le radion participe aux deux lectures : une unité de phase
à la frontière et une unité orientée d'action dans le relèvement.

\paragraph{Rang de mémoire comme lecteur dimensionnel}

Supposons qu'une projection visible conserve la valeur de phase, mais perde
\(d\) cocycles de mémoire linéairement indépendants. Tout relèvement fidèle
nécessite au moins \(d\) coordonnées supplémentaires pour les reconstruire. Cela
motive une lecture précise de la mécanique dimensionnelle :
\begin{equation}
\begin{aligned}
\text{dimension émergente}
&=\text{rang minimal de mémoire indépendante}\\
&\phantom{=}\text{non publiable sur le feuillet visible}.
\end{aligned}
\label{eq:dimension-memory}
\end{equation}
On n'affirme pas que toute dimension physique soit un compteur. On affirme que, dans le
domaine TPK, chaque mémoire indépendante qui doit survivre impose une
coordonnée du relèvement.

\begin{narrativebox}
La dimension apparaît lorsque la projection ne peut plus conserver toute
l'histoire. Le cercle suffit à dire où se trouve la phase ; il ne
suffit pas à dire combien de fois elle est revenue, quel feuillet elle a transporté, quel intervalle
s'est contracté ou quelle retenue a survécu.
\end{narrativebox}

\paragraph{Assemblage de cylindres, de spirales et de trajectoires hélicoïdales}

Les cylindres coinductifs ne sont pas des tubes physiques ajoutés à une spirale. Ce sont
des intervalles emboîtés qui conservent les préfixes et resserrent la valeur. La spirale
décrit la contraction transversale ; l'hélice décrit la mémoire longitudinale ;
le cylindre décrit l'ensemble des valeurs compatibles à une profondeur
finie. À la limite, l'intersection des cylindres sélectionne une section
unique.

L'image spatiale réunit, sans confondre :
\[
\begin{array}{ccl}
\text{angle} &\leftrightarrow& \text{phase visible},\\
\text{hauteur} &\leftrightarrow& \text{retour/mémoire},\\
\text{rayon de cylindre} &\leftrightarrow& \text{incertitude de raffinement},\\
\text{feuillet} &\leftrightarrow& \text{orientation bidirectionnelle}.
\end{array}
\]
Le taux \(3^{-54}\) contracte le cylindre à chaque retour ; il ne s'ajoute pas à l'angle et
ne remplace pas la retenue réelle \(\epsR\).

\subsubsection{Clôtures de phase d'ordres \(2\pi\) et \(4\pi\)}

Dans la projection circulaire, \(2\pi\) restaure la phase. Dans un relèvement orienté
à deux feuillets, un tour peut se fermer sur le feuillet conjugué et nécessiter un
second tour pour restaurer l'orientation :
\[
U_{2\pi}=-I,\qquad U_{4\pi}=I.
\]
La structure coïncide avec la périodicité spinorielle lorsque l'on applique le
foncteur à une représentation \(SU(2)\) ; on n'identifie pas le relèvement TPK à
\(SU(2)\) sans cette application.

La lecture surfacique est compatible :
\[
\mathcal A(2\pi)=\hfull,\qquad
\mathcal A(4\pi)=2\hfull.
\]
Dans un ruban de Möbius, un circuit visible inverse le feuillet et deux circuits
restaurent l'orientation. Le \(2\pi/4\pi\) n'est pas une duplication accidentelle ;
c'est la différence entre retour de phase et retour de l'état enrichi.


\section{Ellipse orientée et variables conjuguées : incertitude minimale et oscillation électromagnétique}
\subsection{Variables symplectiques conjuguées et détermination de l'incertitude minimale}
% Fragmento público generado desde propietarios íntegros.
% Residencia: 52.6.1.
% Fuente: ../incoming/radion_monografia_20260827/manuscrito/sections/07_incertidumbre_oscilador.tex:1-111
\noindent La forme symplectique coordonne action, incertitude et dynamique oscillatoire sans modifier la généalogie de l'état.\par\medskip

\subsubsection{Réalisation de l'état dans un espace de Hilbert}

\paragraph{Matrice de covariance minimale}

Une fois construite l'ellipse d'action, elle peut se réaliser dans un couple
d'opérateurs conjugués
\[
[\widehat Q,\widehat P]=i\hret.
\]
On définit
\begin{equation}
V_\eta=\frac{\hret}{2}
\begin{pmatrix}
e^\eta&0\\[2pt]0&e^{-\eta}
\end{pmatrix}.
\label{eq:covariance}
\end{equation}
Son déterminant est
\begin{equation}
\det V_\eta=\frac{\hret^2}{4}.
\label{eq:cov-det}
\end{equation}
Par conséquent,
\begin{equation}
(\Delta Q)^2=\frac{\hret}{2}e^\eta,
\qquad
(\Delta P)^2=\frac{\hret}{2}e^{-\eta},
\qquad
\Delta Q\,\Delta P=\frac{\hret}{2}.
\label{eq:uncertainty}
\end{equation}
L'inégalité de Robertson--Schrödinger est saturée.

L'annihilateur comprimé
\begin{equation}
\widehat a_\eta=
\frac{e^{-\eta/2}\widehat Q+i e^{\eta/2}\widehat P}
{\sqrt{2\hret}}
\label{eq:a-eta}
\end{equation}
satisfait \([\widehat a_\eta,\widehat a_\eta^\dagger]=1\). Son vide a pour
fonction d'onde
\begin{equation}
\psi_\eta(Q)=
(\pi\hret e^\eta)^{-1/4}
\exp\left(-\frac{Q^2}{2\hret e^\eta}\right).
\label{eq:gaussian}
\end{equation}

\paragraph{Ellipse comme niveau de la fonctionnelle de Mahalanobis}

L’ensemble
\begin{equation}
\mathbb E_{\mathrm{act}}
=\{x\in\R^2:x^{\mathsf T}V_\eta^{-1}x\le4\}
\label{eq:mahalanobis}
\end{equation}
a pour demi-axes
\begin{equation}
a_S=2\Delta Q,\qquad b_S=2\Delta P.
\label{eq:two-sigma}
\end{equation}
Son aire est
\[
4\pi\sqrt{\det V_\eta}=2\pi\hret=\hfull.
\]
L'ellipse HMT construite précédemment coïncide exactement avec le contour à deux
écarts de l'état gaussien minimal.

\begin{proposition}[Chaîne unique du quotient de forme]
Après avoir déclaré les interfaces spectrale, symplectique et de Hilbert,
\begin{equation}
\boxed{
e^{-\eta}
=\sqrt{\frac{A-C^*}{A+C^*}}
=\frac{b_1}{a_1}
=\frac{b_S}{a_S}
=\frac{\Delta P}{\Delta Q}.}
\label{eq:shape-chain-hilbert}
\end{equation}
\end{proposition}

L'ellipse d'action correspond au niveau classique \(H=\hret\omega\) et au
contour \(2\sigma\) du vide comprimé. Elle ne doit pas être confondue avec l'énergie
moyenne du vide, \(\hret\omega/2\). Cette différence d'un facteur deux est un
contrôle physique obligatoire.

\paragraph{Conditions d'invariance de l'ellipse comme objet opératoire}

Dans la formulation conventionnelle, une ellipse d'incertitude peut être considérée comme
un outil graphique pour représenter une matrice de covariance. Ici, l'ordre
causal est inversé : l'ellipse positive existe déjà comme publication de
\((A,C^*,\hret)\) ; l'espace de Hilbert la reconnaît au moyen de
\eqref{eq:covariance}. La réalité ontologique revendiquée ne signifie pas qu'une
particule classique parcourt physiquement un ovale dans le plan \((Q,P)\). Elle signifie
que l'état conserve une forme positive, une action totale et un quotient
d'orientation avant qu'une théorie probabiliste ne les interprète.

La fonction d'onde ne crée pas la forme ; elle la représente. L'incertitude ne crée pas
\(\hbar\) ; elle publie l'impossibilité de comprimer simultanément les deux
directions sans modifier le produit symplectique.

\begin{narrativebox}
La lecture probabiliste est une réalisation possible de l'ellipse, non sa cause.
L'objet premier est la conservation conjointe du produit et de l'orientation. La
probabilité apparaît lorsque cette géométrie est publiée comme covariance d'un
état sur Hilbert.
\end{narrativebox}


\subsection{Échange \(\mu\leftrightarrow\varepsilon\), inversion \(Z\leftrightarrow Z^{-1}\) et conservation de la vitesse causale}
% Fragmento público generado desde propietarios íntegros.
% Residencia: 52.6.2.
% Fuente: ../incoming/radion_monografia_20260827/manuscrito/sections/07_incertidumbre_oscilador.tex:156-197
\paragraph{Échange énergétique entre les secteurs électrique et magnétique}

Les deux termes de \eqref{eq:HLC} deviennent
\begin{equation}
U_E(\theta)=\hret\omega\cos^2\theta,
\qquad
U_B(\theta)=\hret\omega\sin^2\theta.
\label{eq:energy-exchange}
\end{equation}
Par conséquent,
\[
U_E+U_B=\hret\omega.
\]
La polarité électrique--magnétique ne consiste pas à étiqueter arbitrairement
les foyers. Elle consiste en deux réserves conjuguées qui s'échangent pendant
l'orbite tout en maintenant constantes la fréquence et l'action totale.

Le quotient commun se prolonge :
\begin{equation}
\boxed{
e^{-\eta}
=\frac{b_S}{a_S}
=\frac{\Delta P}{\Delta Q}
=\frac{Z_\eta}{Z_*}
=\sqrt{\frac{A-C^*}{A+C^*}}.}
\label{eq:shape-chain-LC}
\end{equation}
Les produits conservés sont
\[
a_Sb_S=2\hret,\qquad
\det V_\eta=\frac{\hret^2}{4},
\qquad
L_\eta C_\eta=\omega^{-2}.
\]

\begin{exactbox}
L'oscillateur LC est la réalisation physique minimale de la double lecture : le
produit conserve l'action et la fréquence ; le quotient conserve l'orientation et
l'impédance. Un cycle complet échange électricité et magnétisme sans
perdre l'aire \(\oint P\,\diff Q=\hfull\).
\end{exactbox}


\subsection{Couple inductif--capacitif \((L_\eta,C_\eta)\) : invariance de \(L_\eta C_\eta=\omega^{-2}\) et orientation de l'impédance}
% Fragmento público generado desde propietarios íntegros.
% Residencia: 52.6.3.
% Fuente: ../incoming/radion_monografia_20260827/manuscrito/sections/07_incertidumbre_oscilador.tex:112-155
\subsubsection{Oscillateur inductif--capacitif associé au couple angulaire}

\paragraph{Construction indépendante des 2 réservoirs conjugués}

Étant donnés \(\omega>0\) et une impédance de référence \(Z_*>0\), nous définissons
\begin{equation}
L_\eta=\frac{Z_*}{\omega}e^{-\eta},
\qquad
C_\eta=\frac{1}{Z_*\omega}e^\eta.
\label{eq:LC-def}
\end{equation}
Alors
\begin{equation}
L_\eta C_\eta=\omega^{-2},
\qquad
\sqrt{\frac{L_\eta}{C_\eta}}=Z_*e^{-\eta}.
\label{eq:LC-invariants}
\end{equation}
La fréquence reste fixe ; l'impédance conserve l'orientation de la
forme.

L'hamiltonien est
\begin{equation}
H_{LC}=\frac{q^2}{2C_\eta}+\frac{\Phi^2}{2L_\eta}.
\label{eq:HLC}
\end{equation}
Avec les variables normalisées
\begin{equation}
Q=\sqrt{Z_*}\,q,\qquad
P=\frac{\Phi}{\sqrt{Z_*}},
\label{eq:LC-map}
\end{equation}
on obtient
\begin{equation}
H_{LC}=\frac\omega2
\left(e^{-\eta}Q^2+e^\eta P^2\right).
\label{eq:HLC-QP}
\end{equation}
Sur \eqref{eq:elipse-param},
\begin{equation}
H_{LC}=\hret\omega.
\label{eq:LC-level}
\end{equation}

% Fuente: ../incoming/radion_monografia_20260827/manuscrito/sections/07_incertidumbre_oscilador.tex:198-239
\paragraph{Oscillateur LC comme unité dynamique élémentaire}

L'oscillateur harmonique traverse presque toute la physique parce qu'il linéarise le
voisinage d'un équilibre, décompose les champs en modes, organise les quanta et
permet de lire la propagation comme une superposition de phases. Dans le radion, cette
universalité reçoit une généalogie concrète : le mode ne commence pas par une
équation différentielle externe, mais par une forme positive dont l'action de tour est
fixe. La dynamique conventionnelle apparaît lorsqu'on choisit une horloge \(\omega\) et un
transducteur \(Z_*\).

La fréquence ne modifie pas l'ellipse d'action : elle change l'énergie du niveau
au moyen de \(E=\hret\omega\). Cette séparation est cruciale. La géométrie fixe
l'action ; l'horloge convertit l'action en énergie.

\subsubsection{Réversion temporelle et propagations conjuguées}

Soit \(H=H^*\) et soit \(J_E\) une structure complexe avec \(J_E^2=-I\).
Supposons un opérateur involutif \(C\) tel que
\[
CJ_EC^{-1}=-J_E,\qquad CHC^{-1}=H.
\]
L'évolution
\begin{equation}
U(t)=\exp\left(-\frac{J_EHt}{\hret}\right)
\label{eq:evolution-J}
\end{equation}
satisfait exactement
\begin{equation}
CU(t)C^{-1}=U(-t).
\label{eq:time-reversal}
\end{equation}
Telle est la forme opératoire minimale de la bidirectionnalité : la conjugaison
échange les prolongations temporellement opposées.

La théorie de l'absorbeur de Wheeler--Feynman et la lecture des antiparticules
comme trajectoires temporellement conjuguées offrent une réalisation historique
ultérieure. L'identité \eqref{eq:time-reversal} est démontrée ; la transformer
en une théorie complète des propagateurs avancés et retardés exige de déclarer
l'action, les conditions de frontière, la causalité et les interactions de l'absorbeur.
La monographie conserve l'intuition physique sans transformer une correspondance
en une égalité de théories.



\section{3 régimes du TRIT : polarités électrique--magnétique et élastique--rotationnelle}
\subsection{Coordination des réalisations elliptique, parabolique et hyperbolique}
% Fragmento público generado desde propietarios íntegros.
% Residencia: 52.7.1.
% Fuente: ../incoming/radion_monografia_20260827/manuscrito/sections/08_complejo_electromagnetismo.tex:1-58
\noindent Les 3 régimes du TRIT admettent des réalisations complexe, électromagnétique et mécanique avec des types différenciés.\par\medskip

\subsubsection{Opérateur quadratique de racine interne de \(-1\)}

\paragraph{Représentation exponentielle du régime elliptique \(J^2=-I\)}

Avec \(J^2=-I\),
\begin{equation}
T_{\mathrm{ell}}(x,y)=e^{-xI-yJ}
=e^{-x}(\cos y\,I-\sin y\,J).
\label{eq:elliptic-exp}
\end{equation}
Sous l'identification générée par \(J\),
\[
T_{\mathrm{ell}}(x,y)\longleftrightarrow e^{-x-iy}.
\]
Le scalaire \(x\) contrôle la dilatation ; \(y\) contrôle la rotation orientée. La forme
\(a+bi\) se reconnaît comme élasticité plus rotation, non comme deux quantités
collées arbitrairement.

\paragraph{Représentation exponentielle du régime parabolique \(J^2=0\)}

Avec \(N^2=0\), l'exponentielle se tronque :
\begin{equation}
T_{\mathrm{par}}(x,y)=e^{-x}(I-yN).
\label{eq:parabolic-exp}
\end{equation}
Le régime décrit le seuil où il n'y a ni rotation périodique ni ouverture
exponentielle. Dans la lecture temporelle active, c'est le présent euclidien, mais
il conserve la mémoire des branches qui le délimitent.

\paragraph{Représentation exponentielle du régime hyperbolique \(J^2=+I\)}

Avec \(R^2=I\) et
\(P_\pm=(I\pm R)/2\),
\begin{equation}
T_{\mathrm{hyp}}(x,y)=e^{-xI-yR}
=q_+P_++q_-P_-,
\qquad q_\pm=e^{-x\mp y}.
\label{eq:hyperbolic-exp}
\end{equation}
Les invariants élémentaires sont
\begin{equation}
q_+q_-=e^{-2x},
\qquad
\frac{q_-}{q_+}=e^{2y}.
\label{eq:q-product-ratio}
\end{equation}
Le produit conserve le mode commun ; le quotient conserve la séparation
orientée.

Appliqué au couple angulaire,
\[
x=\frac{\pi A}{180},\qquad y=\frac{\pi C^*}{180}.
\]
C'est pourquoi \(A\) est la coordonnée commune de la publication et \(C^*\) son
contraste orienté. Ce n'est qu'après l'application d'un lecteur physique qu'ils sont reconnus
comme pôles électriques et magnétiques.


\subsection{Séparation typée des polarités constitutives et mécaniques}
% Fragmento público generado desde propietarios íntegros.
% Residencia: 52.7.2.
% Fuente: ../incoming/radion_monografia_20260827/manuscrito/sections/08_complejo_electromagnetismo.tex:59-158
\subsubsection{Opérateur constitutif sur les canaux \(90/120\)}

\paragraph{Décomposition en composantes commune et orientée}

On définit
\begin{align}
\mathcal E
&=\log\bigl[(1-q_+^{90})(1-q_-^{90})\bigr]
-\log\bigl[(1-q_+^{120})(1-q_-^{120})\bigr],
\label{eq:Ecommon}\\
\mathcal B
&=\log\frac{1-q_-^{90}}{1-q_+^{90}}
-\log\frac{1-q_-^{120}}{1-q_+^{120}}.
\label{eq:Borient}
\end{align}
Sous \(C^*\mapsto-C^*\), \(q_+\leftrightarrow q_-\) s'échangent. Par
conséquent, \(\mathcal E\) est pair et \(\mathcal B\) est impair.

Les lois constitutives normalisées sont
\begin{equation}
\widehat\mu=e^{\mathcal E+\mathcal B},
\qquad
\widehat\varepsilon=e^{\mathcal E-\mathcal B},
\qquad
\widehat Z=e^{\mathcal B},
\qquad
\widehat c=e^{-\mathcal E}.
\label{eq:constitutives}
\end{equation}
On vérifie exactement
\begin{equation}
\widehat Z=\sqrt{\frac{\widehat\mu}{\widehat\varepsilon}},
\qquad
\widehat c=(\widehat\mu\widehat\varepsilon)^{-1/2}.
\label{eq:constitutive-identities}
\end{equation}
La conjugaison produit
\begin{equation}
C^*\mapsto-C^*:\qquad
\widehat\mu\leftrightarrow\widehat\varepsilon,
\quad
\widehat Z\leftrightarrow\widehat Z^{-1},
\quad
\widehat c\mapsto\widehat c.
\label{eq:EM-conjugation}
\end{equation}

\begin{theorem}[Polarité constitutive]
Le lecteur \(90/120\) transforme le mode commun et le contraste orienté du
couple \((A,C^*)\) en deux lois constitutives conjuguées. L'inversion de feuillet
échange perméabilité et permittivité, inverse l'impédance et conserve la
vitesse normalisée.
\end{theorem}

Telle est la formulation précise de l'intuition \enquote{partie électrique et
partie magnétique}. On n'écrit ni \(A=E\) ni \(C^*=B\). On écrit l'application
\[
(A,C^*)\to(x,y)\to(q_+,q_-)
\to(\mathcal E,\mathcal B)
\to(\widehat\varepsilon,\widehat\mu,\widehat Z,\widehat c).
\]

\paragraph{Réalisation élastique--rotationnelle de la polarité constitutive}

La même information peut être publiée dans le régime complexe au moyen de
\begin{equation}
Z_{\mathrm{em}}=e^{\mathcal E/2}
\left(\cos\frac{\mathcal B}{2}\,I
-\sin\frac{\mathcal B}{2}\,J\right).
\label{eq:Zem}
\end{equation}
Le facteur \(e^{\mathcal E/2}\) est une dilatation élastique ; le facteur généré
par \(J\) est une rotation. L'écriture complexe n'ajoute pas une dimension
imaginaire mystérieuse à un monde réel préalablement donné : elle enregistre deux modes
opératoires que le TRIT a déjà séparés.

\paragraph{Transduction de l'orientation en charge effective}

Le feuillet \(\sigma\) du radion fixe une orientation avant d'introduire une
unité de charge. La mécanique dimensionnelle dispose ensuite le transducteur
\[
e_{\mathrm{int}}=\sqrt{\frac{2\alphah\hfull}{Z_{\mathrm{int}}}},
\qquad
Q=n_Qe_{\mathrm{int}}.
\]
La chaîne causale est
\[
\text{feuillet}\to\text{orientation}\to\text{constitutive}
\to\text{transducteur dimensionnel}\to\text{charge physique}.
\]
Appeler l'unité \emph{radion} n'identifie pas prématurément l'orientation
au coulomb ; cela rend visible le lieu où le signe de charge peut résider.

\begin{narrativebox}
L'électricité et le magnétisme ne sont pas collés au cercle depuis l'extérieur. Ils émergent
lorsque le même état est lu à travers le produit et le quotient : le produit
retient le mode commun, le quotient retient l'orientation. L'oscillateur LC
réalise l'échange ; le lecteur \(90/120\) publie les lois constitutives ; la
carte dimensionnelle attribue des unités.
\end{narrativebox}


\section{Séparation des publications \(90/120\) : incidence, registre, spectre, noyau de Barbero et tours}
\subsection{Séparation typée des publications du secteur \(90/120\)}
% Fragmento público generado desde propietarios íntegros.
% Residencia: 52.8.1.
% Fuente: ../incoming/radion_monografia_20260827/manuscrito/sections/09_sector_90120.tex:176-244
\subsubsection{Hiérarchie harmonique globale et semi-groupe \(\langle90,120\rangle\)}

Pour éviter une collision de signes, nous fixons
\[
R_{\mathrm{tow}}:=-R_{\mathrm{ch}}.
\]
Alors
\begin{equation}
T=e^{-xI+yR_{\mathrm{tow}}}
=q_-P_+^{\mathrm{tow}}+q_+P_-^{\mathrm{tow}}.
\label{eq:tower-operator}
\end{equation}
Les tours paires et impaires sont
\begin{align}
c_n&=\operatorname{Tr}(T^n)
=q_-^n+q_+^n
=2e^{-nx}\cosh(ny),
\label{eq:cn}\\
s_n&=\operatorname{Tr}(R_{\mathrm{tow}}T^n)
=q_-^n-q_+^n
=2e^{-nx}\sinh(ny).
\label{eq:sn}
\end{align}
Sous \(C^*\mapsto-C^*\), \(c_n\) est pair et \(s_n\) est impair. On a
\begin{equation}
c_n^2-s_n^2=4e^{-2nx},
\label{eq:tower-invariant}
\end{equation}
et les deux suites satisfont la récurrence d'ordre deux déterminée par les
valeurs propres \(q_\pm\).

Le semi-groupe
\begin{equation}
\langle90,120\rangle=30\langle3,4\rangle
\label{eq:semigroup}
\end{equation}
est la hiérarchie harmonique d'un unique opérateur bicouche. Ce n'est pas une liste de huit
corrections latérales. Sa complétude signifie que tout rapport positif
admet un développement convergent ; elle ne signifie pas qu'un développement obtenu à
partir de la cible soit un sélecteur physique prospectif.

\subsubsection{Classification typée et séparation des réalisations \(90/120\)}

\begin{longtable}{>{\bfseries}p{.19\textwidth} p{.24\textwidth} p{.45\textwidth}}
\toprule
Objet & Opération & Ce qu'il conserve / pourquoi ce n'est pas un autre objet\\
\midrule
\endfirsthead
\toprule
Objet & Opération & Ce qu'il conserve / pourquoi ce n'est pas un autre objet\\
\midrule
\endhead
Incidence \(N_6,N_8\) & dénombrement fini sur les phases actives & cardinalité et orientation ; produit \(90,120\), non une série ;\\
\(-1/12\) & différence normalisée \eqref{eq:minus-one-twelve} & scalaire d'incidence ; non le poids de Barbero ;\\
\(E_{\mathrm{dod}}\) & extracteur harmonique de la roue & phase dodécaphasique ; non la retenue réelle ;\\
\(h=C^*/6\) & reconstruction de six paires & ouverture par paire ; non une normalisation métrologique ;\\
\(F_{\mathrm{spec}}\) & queue logarithmique \(j\ge2\) & publication spectrale complète ; diffère de \(\epsR^\circ\) ;\\
\(\epsone\) & soustraction APP entre \(S_T\) et \(S_E\) & raccord réel de deux relèvements ; non une série de Lambert ;\\
\(K_{6\to8}\) & \(12\Delta S_{90}-\Delta S_{120}\) & saut marqué et évaluation conjuguée ; coefficient de pôle \(1/24\) vérifié ; sortie de Barbero ;\\
\(12:-1\) & douze secteurs et une couture orientée & poids du noyau ; certificat ultérieur \eqref{eq:pole-condition} ; distinct de \(-1/12\) ;\\
\(\langle90,120\rangle\) & semi-groupe des puissances de \(T\) & hiérarchie harmonique globale ; non un résidu par espèce.\\
\bottomrule
\end{longtable}

\begin{thesisbox}
Tous ces objets descendent du même état APP--TRIT--TPK. Cette généalogie
commune explique pourquoi ils partagent des nombres et des symétries. Leur typage explique pourquoi
ils ne peuvent pas se remplacer mutuellement.
\end{thesisbox}

\subsection{Conservation de l'incidence et de l'orientation entre le noyau tritique et ses réalisations ultérieures}
% Fragmento público generado desde propietarios íntegros.
% Residencia: 52.8.2.
% Fuente: ../incoming/radion_monografia_20260827/manuscrito/sections/09_sector_90120.tex:1-30
\noindent La hiérarchie typée sépare incidence, registre, évaluation spectrale, fonctionnelle de Barbero et tours harmoniques.\par\medskip

\subsubsection{Incidence tritique, registre dodécaphasique et lecture sexadique}

\paragraph{Dérivation combinatoire primaire de \(90\) et \(120\)}

Les dénombrements
\[
90=6\binom62,
\qquad
120=8\binom62
\]
naissent du sélecteur de phase et de ses incidences actives. L'orientation
normalisée satisfait
\begin{equation}
\frac{30}{120}-\frac{30}{90}=-\frac1{12}.
\label{eq:minus-one-twelve}
\end{equation}
Ce scalaire est un résultat d'incidence. Ce n'est pas le poids \(12:-1\) du
noyau de Barbero et ce n'est pas une retenue dodécaphasique.

L'extracteur dodécaphasique peut être publié au moyen d'une combinaison
d'harmoniques telle que
\[
E_{\mathrm{dod}}(\phi)=12\cos(3\phi)-\cos(4\phi).
\]
Son domaine est la phase de la roue ; sa sortie est un sélecteur harmonique. Bien que
les nombres \(12\) et \(-1\) apparaissent, il n'opère pas sur des séries de Lambert et ne
produit pas \(\epsR\).


\subsection{Origine dodécaphasique du poids \(12:-1\) et vérification diophantienne de son coefficient de pôle}
% Fragmento público generado desde propietarios íntegros.
% Residencia: 52.8.3.
% Fuente: ../incoming/radion_monografia_20260827/manuscrito/sections/09_sector_90120.tex:115-175
\subsubsection{Fonctionnelle hexade--octade de Barbero--Immirzi}

\paragraph{Séries orientées et saut d'incidence}

On définit
\begin{equation}
S_{90}(q)=\frac{q^{90}}{1-q^{270}},
\qquad
S_{120}(q)=\frac{q^{120}}{1-q^{360}},
\label{eq:barbero-series}
\end{equation}
et
\[
\Delta S_m=S_m(q_-)-S_m(q_+).
\]
La fonctionnelle d'inclusion marquée est
\begin{equation}
K_{6\to8}=12\Delta S_{90}-\Delta S_{120}.
\label{eq:barbero-núcleo}
\end{equation}

\paragraph{Vérification diophantienne du poids dodécaphasique \(12:-1\)}

La chaîne fondamentale dodécaphasique et son lecteur de retour, construits dans
\eqref{eq:l9s102:barbero-cadena-dodecafasica}--%
\eqref{eq:l9s102:barbero-peso-generado}, produisent douze termes de type
\(S_{90}\) et un terme de frontière de type \(-S_{120}\). Le couple orienté
est déterminé avant le calcul du coefficient de pôle.

Considérons \(aS_{90}-bS_{120}\). La vérification du coefficient de
\((1-q)^{-1}\) prend la forme
\begin{equation}
\frac{a}{270}-\frac{b}{360}=\frac1{24}.
\label{eq:pole-condition}
\end{equation}
Multiplier par \(1080\) donne
\[
4a-3b=45.
\]
L'unique couture du tour fondamental a une multiplicité absolue de un.
Sa vérification diophantienne correspond au contre-terme entier minimal
\(|b|=1\). Pour \(b=1\), \(a=12\) ; pour \(b=-1\), \(a=21/2\) n'est pas entier.
Par conséquent,
\begin{equation}
\boxed{a:b=12:1,\qquad aS_{90}-bS_{120}=12S_{90}-S_{120}.}
\label{eq:12minus1}
\end{equation}

\begin{proposition}[Unicité de la vérification diophantienne]
Le couple généré \((12,1)\) est l'unique solution entière de
\eqref{eq:pole-condition} avec canal principal positif et contre-terme
minimal \(|b|=1\).
\end{proposition}

La réalisation hexade--octade conserve les canaux \(90/120\)
préalablement construits ; la chaîne dodécaphasique apporte les multiplicités
et l'orientation de frontière ; le lecteur de retour attribue les séries aux
générateurs. Ces opérations conjointes produisent \(12S_{90}-S_{120}\).
L'équation \eqref{eq:pole-condition}, la positivité et la minimalité
conservent leur fonction de vérification arithmétique ultérieure.

La chaîne de Barbero est ultérieure :
\[
(A,C^*)\to(x,y)\to q_\pm
\to K_{6\to8}\to\gamma_{\mathrm{BI}},
\]
avec l'évaluation
\[
\gamma_{\mathrm{BI}}=0.274007672694459\ldots.
\]
Barbero--Immirzi ne génère pas rétroactivement \(C^*\), \(\Delta_4\) ni
\(\epsR\).


\subsection{Extraction de la composante finie postérieure au système dodécaphasique orienté \(R_{12}\) : \(C^*\mapsto C^*/6\), reste \(j\ge2\) et reconstruction de \(6\) sections}
% Fragmento público generado desde propietarios íntegros.
% Residencia: 52.8.4.
% Fuente: ../incoming/radion_monografia_20260827/manuscrito/sections/09_sector_90120.tex:31-54
\paragraph{Extraction de la composante \(C^*/6\) du registre dodécaphasique}

Les douze secteurs sont regroupés en six paires opposées :
\begin{equation}
p_i=e_i+e_{i+6},\qquad
a_i=e_i-e_{i+6},\qquad i=1,\ldots,6.
\label{eq:sexadic-pairs}
\end{equation}
La reconstruction a pour profil \(1+5+6=12\). L'ouverture par paire est
\begin{equation}
h=\frac{C^*}{6}.
\label{eq:h-C6}
\end{equation}
Ce n'est ni une normalisation extérieure ni le sixième d'une cible : c'est l'indice
interne de la reconstruction sexadique du registre \(\Rstate\).

On définit
\begin{equation}
q_{h,-}=\exp\left[-\frac\pi{180}(A-h)\right],
\qquad
q_{h,+}=\exp\left[-\frac\pi{180}(A+h)\right].
\label{eq:q-h}
\end{equation}


\subsection{Décomposition quantitative des publications de retenue et spectrale : \(F_{\mathrm{spec}}=\varepsilon_R^\circ+\chi_R\)}
% Fragmento público generado desde propietarios íntegros.
% Residencia: 52.8.5.
% Fuente: ../incoming/radion_monografia_20260827/manuscrito/sections/09_sector_90120.tex:55-114
\subsubsection{Publication spectrale de la queue d'ordres \(j\ge2\)}

\paragraph{Définition de la fonctionnelle spectrale}

Pour \(m\in\{90,120\}\), soit
\begin{equation}
L_m^{(2)}(q)=\sum_{j\ge2}\frac{q^{mj}}{j}
=-\log(1-q^m)-q^m.
\label{eq:L2}
\end{equation}
La lecture d'une paire est
\begin{equation}
T_h^{(2)}=\frac{180}{\pi}
\Bigl[
L_{90}^{(2)}(q_{h,-})-L_{90}^{(2)}(q_{h,+})
-L_{120}^{(2)}(q_{h,-})+L_{120}^{(2)}(q_{h,+})
\Bigr].
\label{eq:T-h}
\end{equation}
La reconstruction des six paires donne
\begin{equation}
F_{\mathrm{spec}}=6T_h^{(2)}
=5.1499235153094574410\ldots\times10^{-8}\,{}^\circ.
\label{eq:Fspec}
\end{equation}

La retenue du radion est
\[
F_{\mathrm{acarreo}}=\epsR^\circ
=5.1383016758575282072\ldots\times10^{-8}\,{}^\circ.
\]
Ils ne coïncident pas :
\begin{equation}
\chi_R:=F_{\mathrm{spec}}-F_{\mathrm{acarreo}}
=1.16218394519292338\ldots\times10^{-10}\,{}^\circ.
\label{eq:chiR}
\end{equation}

L’égalité
\[
F_{\mathrm{acarreo}}=F_{\mathrm{spec}}-\chi_R
\]
est un changement de carte valable une fois les deux sorties construites. Il ne doit pas
être présenté comme génération indépendante de \(\epsR^\circ\) si \(\chi_R\) est
défini par la soustraction elle-même. La queue est clôturée et significative ; ce qui est faux,
c'est de l'identifier directement à la retenue.

\begin{falsifier}
Après conversion de la queue exprimée en degrés dans l'unité angulaire interne,
\[
\frac{\pi}{180^\circ}F_{\mathrm{spec}}-\epsone
=2.0283936357433839\ldots\times10^{-12}\ne0.
\]
Par conséquent, ni le log complet ni la queue \(j\ge2\), même après application de la
carte d'unités, ne sont \(\epsone\). Une
seconde factorisation spectrale autonome nécessiterait de générer \(\chi_R\) sans
utiliser la différence comme définition ; la clôture APP--TPK du radion ne dépend pas
de cette promotion.
\end{falsifier}



\section{Holonomie orientée et géométrie physique : réalisations ultérieures de Minkowski, Hodge, Lorentz et Cartan--Holst}
% Fragmento público generado desde propietarios íntegros.
% Residencia: 52.9.
% Fuente: ../incoming/radion_monografia_20260827/manuscrito/sections/10_minkowski_hodge_lorentz.tex:1-55
\noindent L'holonomie orientée admet des réalisations causales et géométriques ultérieures dans Minkowski, Hodge, Lorentz et Cartan--Holst.\par\medskip

\paragraph{Construction de la métrique lorentzienne à partir du bloc APP}

\paragraph{Matrice APP de paramètres \(7\) et \(2\)}

La matrice centrale
\begin{equation}
B_c=\begin{pmatrix}7&2\\2&7\end{pmatrix}
=9P_++5P_-
\label{eq:Bc}
\end{equation}
est normalisée par son déterminant :
\begin{equation}
M_c=\frac{B_c}{\sqrt{45}}
=\exp\left(\frac12\log\frac95\,R\right)
\in SO^+(1,1).
\label{eq:Mc}
\end{equation}
La rapidité interne est
\[
\eta_c=\frac12\log\frac95.
\]
La publication projective associée satisfait
\begin{equation}
P_B^n=P_++\left(\frac59\right)^nP_-.
\label{eq:PB-contract}
\end{equation}
Le même rapport \(5/9\) qui apparaît comme contraction spectrale est reconnu ici
comme séparation d'un mode stable et d'un autre contractant. L'identité est
opératoire ; elle ne transforme pas toutes les occurrences de \(5/9\) en la même application.

\paragraph{Incidence \(6\to8\) et quotient de signature causale}

Soit \(M\) l'incidence typée entre six faces et huit sommets. La relation
spectrale récupérée est
\begin{equation}
MM^{\mathsf T}=12P_u+4P_-,
\label{eq:MMt}
\end{equation}
où \(P_u\) projette sur le mode uniforme et \(P_-\) sur le sous-espace
orienté. Le quotient survivant de dimension quatre se décompose comme
\begin{equation}
Q_4\simeq\R u\oplus E_-.
\label{eq:radion-Q4}
\end{equation}
En déclarant négatif le mode uniforme et positifs les trois modes orientés,
on obtient
\begin{equation}
B_{\mathrm{caus}}=\operatorname{diag}(-1,1,1,1).
\label{eq:minkowski}
\end{equation}
L'écran de Minkowski ne sélectionne pas l'état HMT. C'est la réalisation
causale ultérieure du quotient que l'incidence a déjà construit.


\subsection{Dualité de Hodge sur les \(2\)-formes lorentziennes}
% Fragmento público generado desde propietarios íntegros.
% Residencia: 52.9.1.
% Fuente: ../incoming/radion_monografia_20260827/manuscrito/sections/10_minkowski_hodge_lorentz.tex:56-100
\subsubsection{Dualité de Hodge et opérateur de quart de tour}

La signature \((-+++ )\) et une orientation étant fixées, l'étoile de Hodge sur
les deux-formes satisfait
\begin{equation}
\star:\Lambda^2Q_4^*\to\Lambda^2Q_4^*,
\qquad \star^2=-I.
\label{eq:hodge-star}
\end{equation}
De cette manière, le quart de tour qui, dans le plan de phase, s'exprimait au moyen de
\(J^2=-I\) réapparaît dans l'espace des champs comme structure complexe sur
\(\Lambda^2\). Les domaines sont différents, mais la généalogie d'orientation
est conservée par le foncteur.

Soit \(A_{\mathrm{em}}\) un potentiel et
\begin{equation}
F=\diff A_{\mathrm{em}}.
\label{eq:F-dA}
\end{equation}
La décomposition de \(F\) par rapport à un observateur temporel produit des champs
électrique et magnétique ; \(\star F\) échange leurs rôles avec le signe fixé
par l'orientation et la signature. La polarité constitutive de
\eqref{eq:EM-conjugation} acquiert ainsi une réalisation différentielle.

\paragraph{Chaîne de transduction électromagnétique et causale}

Le transport qui doit rester visible est
\begin{equation}
\Gamma_9
\longrightarrow\text{hélice de mémoire}
\longrightarrow(Q_4,B_{\mathrm{caus}},\star)
\longrightarrow A_{\mathrm{em}}
\longrightarrow F=\diff A_{\mathrm{em}}
\longrightarrow
m\nabla_u u=q(\iota_uF)^\sharp.
\label{eq:lorentz-chain}
\end{equation}
Chaque flèche ajoute de la structure : l'holonomie conserve le retour ; le quotient fixe
la causalité ; Hodge fixe la dualité ; le potentiel fixe le champ ; le couplage
avec la charge produit la dynamique.

L'équation finale est la force de Lorentz sous forme géométrique. L'hélice TPK
est son antécédent de mémoire, non une orbite physique déjà mue par un champ. La
dynamique n'apparaît qu'à l'introduction de \(F\), \(q\) et de la connexion \(\nabla\).


\subsection{Séparation entre mémoire torsionnelle interne et réalisation géométrique de Cartan--Holst}
% Fragmento público generado desde propietarios íntegros.
% Residencia: 52.9.2.
% Fuente: ../incoming/radion_monografia_20260827/manuscrito/sections/10_minkowski_hodge_lorentz.tex:101-186
\subsubsection{Réalisation lorentzienne des trajectoires hélicoïdales}

Dans un champ magnétique uniforme, une particule de charge \(q\), de masse \(m\) et de
facteur relativiste \(\gamma\) possède
\begin{equation}
\omega_c=\frac{|q|B}{\gamma m},
\qquad
r_L=\frac{p_\perp}{|q|B},
\qquad
\Delta z=\frac{2\pi p_\parallel}{|q|B}.
\label{eq:lorentz-helix}
\end{equation}
La phase cyclotronique se ferme tandis que le mouvement parallèle avance : l'orbite
est hélicoïdale. Le parallélisme avec \eqref{eq:H9-helix} est structurel et peut
être formalisé par un transducteur qui préserve phase, orientation et pas. On
n'identifie pas le quotient \(q_9\) à la coordonnée spatiale \(z\) ; on reconnaît
la même architecture de retour plus avance.

La force de Lorentz ne se dérive pas de l'image d'une spirale par ressemblance.
Elle se dérive en complétant \eqref{eq:lorentz-chain}. L'image aide à comprendre
pourquoi phase circulaire et mémoire longitudinale sont compatibles ; la forme de
champ détermine la dynamique.

\subsubsection{Confluence entre action, aire, information et gravité}

\paragraph{Réalisation Cartan--Holst de l'holonomie}

La promotion gravitationnelle s'exprime au moyen d'un morphisme d'holonomies
\begin{equation}
\operatorname{Hol}_{\mathcal A}
\bigl(\mathfrak R_{\mathrm{grav}}(\gamma)\bigr)
=\rho_C\bigl(\operatorname{Hol}_{\TPK}(\gamma)\bigr).
\label{eq:holonomy-gravity}
\end{equation}
Le membre de droite contient la mémoire de chemin ; \(\rho_C\) la réalise dans une
connexion gravitationnelle. La torsion de Cartan
\[
T^I=D_\omega e^I
\]
apparaît dans le codomaine. La torsion globale \(\Delta_4\) est une coordonnée
en amont qui peut alimenter l'application, mais n'est pas renommée \(T^I\) sans
\eqref{eq:holonomy-gravity}.

\paragraph{Transduction de l'aire d'action en aire informationnelle}

Le radion fixe un quantum d'action par phase. Le lecteur hexade--octade de
Barbero utilise ultérieurement les canaux \(q_\pm\), l'incidence \(90/120\) et
la fonctionnelle \eqref{eq:barbero-núcleo}. La direction causale est
\[
(A,C^*,\hret)\to q_\pm\to K_{6\to8}
\to\gamma_{\mathrm{BI}}\to\text{aire/information}.
\]
L'alphabet tritique apporte l'unité d'entropie \(\log3\) et une coupure de
\(81\) canaux apporte \(81\log3\). Cette entropie ne se confond ni avec
\(\log2\), ni avec \(\log\varphi\), ni avec une entropie microcanonique, ni avec une
entropie de von Neumann sans spécifier l'état.

La synthèse physique est affirmative mais typée : le radion apporte l'unité
d'action orientée qu'une holonomie peut transporter ; Barbero publie la
normalisation d'aire ; l'écran causal et Hodge permettent de réaliser des champs ;
Cartan--Holst reçoit l'holonomie en géométrie gravitationnelle.

\begin{statusbox}
L'équation scalaire du radion n'est pas à elle seule une équation d'Einstein. Sa
contribution à la gravité consiste à être une section exacte de l'architecture
d'action, d'orientation et d'holonomie que la réalisation gravitationnelle consomme.
\end{statusbox}

\subsubsection{Transformation de Wick comme changement de lecteur analytique}

Dans la formulation conventionnelle, la rotation de Wick relie une variable
temporelle lorentzienne à une coordonnée euclidienne par un prolongement
complexe. Dans HMT, la racine \(J^2=-I\), l'opérateur involutif \(R^2=I\) et le seuil
\(N^2=0\) existent déjà auparavant. Le changement de lecteur
\[
J\longleftrightarrow R
\]
transporte un couple commun/orienté entre publication elliptique et
publication hyperbolique. Le prolongement complexe reconnaît ensuite cette
opération en coordonnées analytiques.

Cela explique pourquoi cercle, ellipse d'action et intérieur hyperbolique peuvent
appartenir au même état sans être la même métrique. Le cercle utilise \(J\) pour
la phase ; la forme de l'ellipse utilise une compression symplectique générée par \(R\) ; Klein mesure
l'intérieur par le birapport ; l'écran de Minkowski réalise la causalité
dans un autre codomaine.


\section{Ellipse, réseau et phase de chemin : caractère modulaire, orbites et propagation}
\subsection{Réseau rectangulaire de périodes, involution \(\tau\mapsto-1/\tau\) et invariance de \(j\)}
% Fragmento público generado desde propietarios íntegros.
% Residencia: 52.10.1.
% Fuente: ../incoming/radion_monografia_20260827/manuscrito/sections/11_modularidad_orbitas.tex:1-84
\noindent Les caractères modulaires, orbitaux et de chemin conservent des invariants distincts d'une même forme orientée.\par\medskip

\subsubsection{Tore complexe associé à l'ellipse d'action}

\paragraph{Réseau rectangulaire de périodes}

L'ellipse à elle seule n'est pas un tore. Nous déclarons le réseau
\begin{equation}
\Lambda_\eta=a_S\Z+i,b_S\Z
\label{eq:lattice}
\end{equation}
et le quotient \(\C/\Lambda_\eta\). Son module est
\begin{equation}
\tau_\eta=i\frac{b_S}{a_S}=ie^{-\eta}
=i\,0.741048144159375160795483663369300\ldots.
\label{eq:tau-eta}
\end{equation}
L'inversion bidirectionnelle \(\eta\mapsto-\eta\) échange les axes :
\begin{equation}
\tau_{-\eta}=ie^\eta=-\frac1{\tau_\eta}.
\label{eq:modular-S}
\end{equation}
C'est exactement la transformation modulaire \(S:\tau\mapsto-1/\tau\).

\paragraph{Invariance de la fonction modulaire \(j\)}

Par invariance modulaire,
\begin{equation}
\boxed{j(\tau_{-\eta})=j(\tau_\eta).}
\label{eq:j-invariance}
\end{equation}
Dans la normalisation \(j(i)=1728\),
\[
j(\tau_\eta)
=5597.43843932408729830097089254768\ldots.
\]
La fonction moonshine normalisée \(J=j-744\) donne
\[
J(\tau_\eta)
=4853.43843932408729830097089254768\ldots.
\]
Pour le bloc APP basal,
\[
\tau_0=i\frac{\sqrt5}{3},
\qquad
j(\tau_0)=5369.48832024674306021916882626352\ldots.
\]

\begin{proposition}[Ce que conserve la classe modulaire]
L'involution de feuillet inverse \(\eta\), mais \(j\) conserve la classe du
tore. Les foyers et l'étiquetage des axes retiennent l'orientation que \(j\)
oublie.
\end{proposition}

Ce résultat relie de manière exacte la double voie à la modularité une
fois le quotient déclaré. Il n'autorise pas à faire agir directement le
Monstre sur la charge, les chiffres ou les foyers. Le couloir exceptionnel
\[
\mathrm{APP}\to\mathrm{Paley}\to\mathrm{Witt}\to\mathrm{Golay}
\to A_2^{12}\to\mathrm{Leech}\to V^\natural
\to\mathrm{Monster}\to\mathrm{Moonshine}
\]
est une autre projection du même état enrichi. Pour identifier une action
concrète sur le radion, il faut un entrelaceur qui préserve ses
coordonnées pertinentes.

\paragraph{Chaîne d'invariants du caractère de forme}

En réunissant les réalisations déclarées,
\begin{equation}
\boxed{
r_\eta=e^{-\eta}
=\sqrt{\frac{A-C^*}{A+C^*}}
=\frac{b_1}{a_1}
=\frac{b_S}{a_S}
=\frac{\Delta P}{\Delta Q}
=\frac{Z_\eta}{Z_*}
=\Im\tau_\eta.}
\label{eq:shape-master}
\end{equation}
Un seul rapport parcourt géométrie spectrale, action, covariance, impédance et
module. Chaque égalité repose sur une application explicite ; aucune ne dépend d'une
ressemblance décimale.


\subsection{Compatibilité entre le caractère modulaire et la propagation sur des chemins orientés}
% Fragmento público generado desde propietarios íntegros.
% Residencia: 52.10.2.
% Fuente: ../incoming/radion_monografia_20260827/manuscrito/sections/11_modularidad_orbitas.tex:85-187
\subsubsection{Conservation surfacique dans les orbites de Kepler et de Sommerfeld}

\paragraph{Équivalence surfacique des réalisations orbitale et oscillatoire}

L'ellipse de Kepler vit dans l'espace de configuration et répond à un potentiel
central. L'ellipse du radion vit dans l'espace des phases et répond à une forme
d'action. Ce sont des objets différents. Elles partagent une loi surfacique parce que les deux
réalisations possèdent une deux-forme conservée.

Pour une orbite képlérienne,
\begin{equation}
r(\phi)=\frac{p}{1+e\cos(\phi-\phi_0)},
\qquad
\frac{\diff A}{\diff t}=\frac{\ell}{2m}.
\label{eq:kepler}
\end{equation}
La troisième loi prend la forme
\[
T^2=\frac{4\pi^2\mu}{K}a^3.
\]
Pour l'oscillateur en phase,
\begin{equation}
J=\oint p\,\diff q=\frac{2\pi E}{\omega}.
\label{eq:osc-action}
\end{equation}
Si la cellule élémentaire est \(J=\hfull\), alors
\begin{equation}
E=\hret\omega,\qquad \frac{J}{2\pi}=\hret.
\label{eq:E-hbarw}
\end{equation}
Telle est la convergence nette entre action de Planck, fréquence angulaire et
radion.

Sommerfeld quantifie les intégrales d'action sur les cycles. La lecture HMT ajoute
la mémoire du cycle, du feuillet et du retour qui survivent. Dans un ruban de Möbius,
un lacet visible peut porter \(\hfull/2\) et deux lacets restaurer \(\hfull\).
Ce facteur topologique ne s'ajoute pas à l'indice de Maslov d'une quantification EBK ;
ils relèvent de corrections ayant des propriétaires distincts.

\paragraph{Holonomie orbitale et propagations avancée et retardée}

Une orbite qui accumule de l'holonomie satisfait une condition du type
\begin{equation}
\epsilon_\gamma\exp\left(\frac{iJ_\gamma}{\hret}\right)=1.
\label{eq:holonomy-action}
\end{equation}
Le signe \(\epsilon_\gamma\) conserve le feuillet ; \(J_\gamma\) conserve
l'action. L'avance ou le retard apsidal s'interprète comme découplage entre
la clôture angulaire visible et le retour complet. Le radion fournit
l'unité orientée avec laquelle cette différence est mesurée.

\subsubsection{Composition des amplitudes de Schrödinger et somme des chemins}

\paragraph{Quatre réalisations d'une même phase}

Une fois \(\hret\) fixé, différentes équations conventionnelles reconnaissent
des aspects du même transport :
\begin{enumerate}
\item Schrödinger publie la phase au moyen de
\(i\hret\partial_t\psi=H\psi\).
\item Pauli ajoute le double feuillet de spin et son couplage magnétique.
\item Dirac linéarise la relation énergie--quantité de mouvement et organise particule et
antiparticule dans une représentation causale.
\item La somme des chemins attribue à chaque histoire la phase
\(\exp(iS[\gamma]/\hret)\).
\end{enumerate}
Ces théories ne génèrent pas rétrospectivement \(\hret\). Elles reçoivent l'unité
d'action et réalisent différents contenus de l'état enrichi.

\paragraph{Fonctionnelle de phase sur l'espace des chemins}

Dans la formulation de Feynman,
\begin{equation}
\mathcal K(b,a)=\int_{\gamma:a\to b}
\exp\left(\frac{i}{\hret}S[\gamma]\right)\mathcal D\gamma.
\label{eq:path-integral}
\end{equation}
Un incrément d'action \(\hret\) change la phase d'un radian. Un incrément
\(\hfull\) la change d'un tour. Le théorème
\(\mathcal A(\theta)=\hret\theta\) transforme cette relation en géométrie
d'aire : la phase d'un chemin peut se lire comme action balayée en unités de
radion.

L'inversion \(C_M=-\operatorname{rev}\) sur les mots dodécaphasiques et
l'identité \eqref{eq:time-reversal} organisent des chemins conjugués. Dans une
réalisation d'absorbeur, on sépare
\[
G_{\mathrm{sym}}=\tfrac12(G_{\mathrm{ret}}+G_{\mathrm{adv}}),
\qquad
G_{\mathrm{odd}}=G_{\mathrm{ret}}-G_{\mathrm{adv}}.
\]
Le premier secteur est pair sous l'échange retardé--avancé ; le second est
impair. Cela reproduit la séparation entre énergie/forme paire et action
orientée impaire du radion lorsque le transducteur préserve la frontière et la forme
symplectique.

\begin{statusbox}
La théorie classique de l'absorbeur et l'interprétation de Feynman--Stückelberg des
antiparticules sont des réalisations historiques différentes. La première utilise une
interaction temporellement symétrique ; la seconde réorganise la propagation dans la
théorie relativiste. HMT apporte un antécédent commun de réversion, mais ne les
fusionne pas sans une application de propagateurs.
\end{statusbox}


\section{Radical nonadique et frontière décimale : conservation de la retenue du radion}
\subsection{La retenue du radion dans la complétion décimale}
% Fragmento público generado desde propietarios íntegros.
% Residencia: 52.11.1.
% Fuente: ../incoming/radion_monografia_20260827/manuscrito/sections/12_aritmetica_borde.tex:1-159
\noindent La frontière nonadique conserve le défaut somme--produit, la vacance et la retenue de complétion.\par\medskip

\subsubsection{Radical \(3\! -\! 6\! -\! 9\) du bord}

\paragraph{Génération et nilpotence du radical nonadique}

Dans l'anneau \(\Z/9\Z\), l'idéal
\begin{equation}
\mathfrak m=(3)=\{0,3,6\}
\label{eq:radical-369}
\end{equation}
satisfait
\begin{equation}
\mathfrak m^2=(9)=0.
\label{eq:m-square-zero}
\end{equation}
C'est le secteur nilpotent du bord nonadique. L'orbite visible de ses trois
marques se publie comme
\[
369\longrightarrow693\longrightarrow936\longrightarrow369.
\]
Le retour n'implique pas que tous les états soient égaux : la permutation
visible se ferme tandis que les cocycles de chemin peuvent avancer.

La double convention
\[
9\equiv0\pmod9,
\qquad
\operatorname{dr}_9(9)=9
\]
exprime deux lecteurs : zéro résiduel algébrique et clôture visible de la racine
numérique. C'est pourquoi un neuf peut absorber une torsion de bord sans devenir
un zéro spécifique de la fonction zêta.

\paragraph{Défaut somme--produit \(459-405=54\)}

Les feuillets APP produisent des totaux complémentaires
\[
S_+=405,\qquad S_\times=459,\qquad S_\times-S_+=54.
\]
Le \(54\) est un défaut exact des tables, et réapparaît comme exposant de
contraction par l'identité \(9\cdot6=54\). Cependant, chaque occurrence
nécessite son application. La distance focale \(d_1=0.544545\ldots\) n'est pas égale à
\(0.54\), et l'analyse des invariants n'a pas produit de flèche qui
l'identifie au défaut \(54\).

La puissance
\[
9^9=387,420,489
\]
génère une échelle où
\begin{equation}
\operatorname{ord}_{10}\bigl(9^{9^9}\bigr)=369,693,099,
\qquad
\#\,\text{chiffres}=369,693,100.
\label{eq:369693}
\end{equation}
Le bloc \(369|693\) apparaît dans le couple ordre--longueur ; on n'affirme pas qu'il soit
le préfixe décimal de la puissance.

\subsubsection{Théorie arithmétique des vacances décimales}

\paragraph{Pente asymptotique de vacance}

La séparation entre la décennie et la base nonadique est
\begin{equation}
\delta_{\mathrm{vac}}=\log_{10}\frac{10}{9}.
\label{eq:delta-vac}
\end{equation}
Le compteur cumulé
\begin{equation}
V_9(n)=\left\lceil n\delta_{\mathrm{vac}}\right\rceil
\label{eq:vacancy-count}
\end{equation}
enregistre le nombre de positions de correction nécessaires à la publication décimale
d'une évolution nonadique. Ses sauts forment un mot binaire
\[
z_n=V_9(n)-V_9(n-1)\in\{0,1\}.
\]
Les lacunes se répartissent avec des longueurs \(21/22\), les retours avec des
longueurs \(6/7\), et le déterminant combiné est \(15\). Ces nombres
proviennent du lecteur de vacance, non du noyau de Barbero, même si le \(15\)
coïncide avec \(\binom62\).

\paragraph{Polarisation et courant arithmétique du réseau nonadique}

On définit
\begin{equation}
P_N=\sum_{n=1}^N(z_n-\delta_{\mathrm{vac}}),
\qquad
J_N=P_N-P_{N-1}=z_N-\delta_{\mathrm{vac}}.
\label{eq:arith-current}
\end{equation}
La séquence équilibrée satisfait \(|P_N|<1\). Le \enquote{courant}
\(J_N\) est ici un courant arithmétique : il mesure la création et l'absorption de
vacances par rapport à la pente moyenne. Seul un transducteur dimensionnel
ultérieur peut le publier comme densité de courant électromagnétique.

L'intuition physique est puissante parce que la même forme algébrique sépare une
moyenne paire et une fluctuation orientée. La rigueur exige de conserver le type :
\[
J_N\in\R\quad\text{sans dimension},
\qquad
j^\mu\in\text{droite physique de courant}.
\]

\subsubsection{Complétion \(729\to1000\) et conservation de la retenue}

\paragraph{Décomposition \(1000=729+243+27+1\)}

La complétion
\[
1000=729+243+27+1
\]
conserve l'unité au moyen de quatre échelles ternaires. La couronne totale est
\[
1000-729=271,
\]
tandis que le dernier entier avant la clôture satisfait
\[
999=729+270.
\]
La différence entre \(270\) et \(271\) distingue frontière occupée et unité
complétante. Le même soin évite d'identifier une vacance à un zéro ou une
retenue à un résidu réel.

\paragraph{Résidence de la retenue du radion dans la couronne décimale}

Les opérandes du radion commencent, en blocs de trois chiffres, par
\[
\Phi_T:\quad000|027|455|906|837|565|\cdots,
\]
\[
D_E:\quad000|027|455|010|034|743|\cdots.
\]
La soustraction APP donne
\[
\epsone:\quad000|000|000|896|802|822|\cdots.
\]
La petitesse du résultat ne s'obtient pas en tronquant les termes inférieurs.
Elle s'obtient parce que deux publications distinctes partagent un long préfixe et que
l'opérateur \(\operatorname{SubAcarreo}_{1000}\) conserve les emprunts nécessaires
pour les soustraire exactement.

Le mot \emph{mémoire} a ici deux niveaux :
\begin{enumerate}
\item \(q_{\mathrm{mem}}\) enregistre la profondeur des retours nonadiques ;
\item \(c_i\) enregistre les emprunts locaux de la soustraction décimale.
\end{enumerate}
Le premier organise l'hélice ; le second publie la décimale. \(\epsR\) utilise
les deux parce que ses opérandes proviennent de cylindres en profondeur et que sa soustraction
s'effectue avec retenue, mais il n'est égal à aucun des compteurs.

\begin{thesisbox}
Le bord n'est pas un lieu où la théorie perd de l'information. C'est le lieu où
l'information change de représentation : phase qui revient, mémoire qui
s'élève, cylindre qui se contracte, vacance qui s'enregistre et retenue qui
transporte la différence.
\end{thesisbox}

\subsection{Compatibilité des sections directe et torsionnelle : caractérisation structurelle de \(\pi\)}
% Fragmento público generado desde propietarios íntegros.
% Residencia: 52.11.2.
% Fuente: ../incoming/radion_monografia_20260827/manuscrito/sections/03_cierre_exacto.tex:286-309
\subsubsection{Caractérisation structurelle de \(\pi\) par compatibilité des sections}

Réordonner \eqref{eq:cierre-unitario} donne
\begin{equation}
\boxed{
2\pih=
\frac{1+\Phi_T-\epsone}
{\frac5{36}+\frac{25}{9}\alphah}.}
\label{eq:pi-radion}
\end{equation}
Cette égalité peut être appelée une nouvelle définition de \(\pi\) au sein de la
carte du radion si l'on entend \emph{définition structurelle} comme une
caractérisation exacte par compatibilité des sections. Ce n'est pas un second
générateur indépendant : \(\Delta_4\) et l'opérateur de retenue reçoivent déjà la
section coinductive de \(\pi\). La nouveauté réside dans le fait que le tour est
caractérisé par le raccord entre publication directe, torsion et mémoire.

\begin{narrativebox}
L'équation \eqref{eq:pi-radion} ne dit pas que nous ignorons d'abord \(\pi\) et
le devinons en fermant un angle. Elle dit que le \(\pi\) généré par la connexion
nonadique est exactement celui qui rend compatibles deux lectures internes du
même état. C'est une identité de reconnaissance intrinsèque, non un
étalonnage métrologique.
\end{narrativebox}


\section{Conservation évolutive de l'unité holographique à travers les dimensions}
% Conversión TeX fiel del cuerpo científico del delta narrativo.
% Fuente congelada: ../../../../../HMT2/CONTINUIDAD_EDITORIAL/RECONSTRUCCION_ARMONICA_MACROTRATADO_20260822/REVISION_CONJUNTA_PARTES_I_II/auditar_pdf_rigor_academico/docs/DELTA_NARRATIVO_UNIDAD_RADION_AUTOCONSERVACION_AUTOINERCIA_GRAVEDAD_PENTADICA_20260828.md:12-366.
\subsection{Séparation typée entre le régime de courbure \(\tau\in\{+1,0,-1\}\), la réponse radiale \(p\) et l'holonomie globale de spin}

La biographie du radion peut être représentée par une courbe relevée

\[
\gamma(\theta)
=
\bigl(r(\theta)\cos\theta,\,
r(\theta)\sin\theta,\,
m(\theta)\bigr),
\]

où les deux premières coordonnées appartiennent au plan euclidien \(\mathbb R^2\) et \(m(\theta)\) enregistre la mémoire ou la profondeur du relèvement. Cette équation permet de lire conjointement cercle, hélice et spirale.

Lorsque l'angle évolue et que le rayon et la mémoire restent constants, la publication plane est un cercle. Lorsque l'angle et la mémoire évoluent tandis que le rayon conserve sa valeur, la trajectoire complète est une hélice de rayon constant. Lorsque l'angle et le rayon évoluent sur une mémoire fixe, une spirale apparaît dans \(\mathbb R^2\). Lorsque les trois coordonnées évoluent, la trajectoire est une hélice radiale : une biographie dans laquelle phase, amplitude et mémoire se transforment conjointement.

Le caractère radial \(p\) classe la loi de réponse du rayon. Une famille utile est

\[
r_p(\theta)
=
r_0\bigl[1+p\lambda(\theta-\theta_0)\bigr]^{1/p},
\qquad p\neq 0,
\]

avec pour limite logarithmique

\[
r_0\exp\bigl(\lambda(\theta-\theta_0)\bigr)
\qquad\text{lorsque }p\to0.
\]

Dans cette paramétrisation, \(p=2\) publie le régime de Fermat, \(p=1\) le régime archimédien, \(p=0\) le logarithmique, \(p=-1\) l'hyperbolique et \(p=-2\) celui du lituus, avec les choix de domaine et d'orientation correspondants. Les valeurs fractionnaires décrivent des régimes intermédiaires. Ainsi, \(p=\tfrac12\), \(p=\tfrac32\) ou \(p=-\tfrac12\) expriment des réponses radiales qui interpolent entre les familles entières et élargissent la classification géométrique disponible.

Le caractère peut aussi changer lors du franchissement d'un seuil. L'écriture différentielle

\[
\frac{dx}{d\theta}
=
\beta(m)\,x^{1-p(m)},
\qquad x=\frac r{r_0},
\]

décrit une biographie par régimes : \(p(m)\) reste stable au sein d'une phase et change lorsque la mémoire atteint une frontière de publication. La trajectoire résultante relie des secteurs circulaires, hélicoïdaux et spiraux par une loi commune. La continuité de l'état et le registre du seuil font de ces secteurs les phases d'une même évolution.

Trois invariants occupent ici des niveaux distincts et coordonnés. Le signe de courbure

\[
\tau\in\{+1,0,-1\}
\]

classe les régimes elliptique, parabolique et hyperbolique. Le caractère \(p\) classe la réponse radiale. Le spin enregistre l'holonomie d'orientation accumulée par le relèvement. Courbure, réponse radiale et spin forment un système de lecture : la courbure détermine le type local d'espace, \(p\) détermine comment la trajectoire occupe cet espace et le spin conserve l'orientation acquise en le parcourant. Les courbes fermées et ouvertes apparaissent alors comme des biographies géométriques distinctes de la même unité orientée.

\subsection{Particule comme section persistante d'une orbite enrichie : phase, chemin, fibre, multisection et mémoire}

La mécanique dimensionnelle peut formuler la particule comme une section persistante de cette extension géométrique et dynamique. Une notation conceptuelle compacte est

\[
\mathfrak P
=
\operatorname{Sec}_{\mathrm{pers}}
\bigl[
\tau,p,\operatorname{Hol};
\text{feuillet},\text{fibre},\text{route},\text{mémoire}
\bigr].
\]

La particule conserve un chemin reconnaissable à travers les changements de phase, de courbure et de dimension. Son spin publie l'holonomie d'orientation ; sa charge publie la polarité de phase ; sa couleur publie une orientation interne de la fibre de jauge ; sa saveur distingue des familles de chemins ou de multisections ; sa masse exprime le caractère de transduction dimensionnelle par lequel l'action surfacique acquiert une lecture volumique. Chaque propriété est un lecteur de l'état persistant et toutes renvoient à une généalogie commune.

L'électron occupe le seuil basal de cette construction parce qu'il offre le premier lecteur stable dans lequel phase, charge, spin et masse sont simultanément publiés. La masse électronique directrice

\[
m_e^\beta=0.510998950690000808608\ldots\ \mathrm{MeV}
\]

fixe une référence de transduction ; les états basaux antérieurs conservent leur fonction d'étapes généalogiques. À partir de l'électron, les quotients \(K_D/K_\Omega\), les tours \(90/120\) et le domaine \(81/56/13/324\) organisent la publication des familles massives. La loi des masses lit donc comment chaque section persistante traverse l'espace enrichi et convertit la densité surfacique d'action en présence volumique.

Cette interprétation relie les spirales aux particules de manière précise. L'indice \(p\) apporte la réponse radiale ; \(\tau\) apporte le régime de courbure ; l'holonomie apporte le spin ; le chemin et la fibre apportent la saveur et la couleur ; la transduction aire–volume apporte le caractère massif. L'atlas des particules devient ainsi une classification des modes stables de traversée dimensionnelle.

\subsection{Transduction de la phase orientée en polarisation électrique, réponse magnétique et émission radiative de frontière}

L'électricité se présente comme transport dirigé de phase et d'action à travers une structure de vacances. Une vacance est une transition admissible du support : elle ouvre un canal, fixe un seuil et permet à l'actualisation d'avancer. La composante électrique publie le transport tangentiel le long du chemin. La composante magnétique publie la courbure transversale induite par ce transport. Leur perpendicularité exprime la relation géométrique entre avance et rotation.

Lorsque deux composantes transversales sinusoïdales maintiennent le déphasage approprié, le vecteur du champ tourne. La propagation longitudinale relève cette rotation en une hélice. La variation d'amplitude transforme sa projection frontale en une spirale. La lecture géométrique est alors entièrement coordonnée : l'angle porte la phase ; la vitesse angulaire porte la fréquence ; le rayon porte l'amplitude et la densité d'action ; l'orientation porte la polarisation et le sens ; la coordonnée verticale porte la propagation et la mémoire ; \(p\) porte le régime radial ; le seuil porte l'émission, l'absorption ou le changement d'état.

L'électromagnétisme apparaît comme une dynamique de phase–action–mémoire avec des publications complémentaires. L'électricité enregistre l'avance polaire ; le magnétisme enregistre la réponse transversale ; l'onde enregistre la périodicité ; l'hélice enregistre l'avance avec mémoire ; la spirale enregistre l'évolution radiale ; le rayonnement enregistre la publication de l'état lorsque l'accumulation atteint un seuil.

Cette architecture ordonne aussi le rôle des vacances et des cycles. La vacance fournit la possibilité locale de transition ; le cycle conserve la phase ; la mémoire distingue les tours ; le seuil décide de la résidence suivante de l'action. La propagation peut rester confinée, changer de régime radial ou se publier comme rayonnement contrôlé. Les forces de champ acquièrent ainsi une description généalogique continue, du support discret jusqu'à l'onde observée.

\subsection{Réversibilité de la dynamique complète et coût thermodynamique des projections réductrices d'information}
Le Landauer nul constitue un cas limite de cette loi. Pour une actualisation bijective de l'état complet,

\[
H(X\mid U(X))=0,
\qquad
\inf Q_L=0.
\]

L'information reste disponible dans la généalogie. Lorsqu'un lecteur réalise un quotient \(q\), l'entropie se décompose comme

\[
H(X)=H(q(X))+H(X\mid q(X)).
\]

Le terme conditionnel mesure l'information retenue hors de la projection. Une réinitialisation tritique matérialisée satisfait le seuil

\[
Q_{\mathrm{reset}}\ge k_B\Theta\log 3.
\]

La constante de Boltzmann traduit la multiplicité que le lecteur thermique regroupe ; le seuil de Nyquist sépare l'échantillonnage fidèle du repliement spectral ; le Landauer nul caractérise le transport réversible de l'état intégral. Information, température et action sont articulées par le type de lecture appliqué.

\subsection{Forces autoconservatives comme transport d'action entre les secteurs observable, radial, de mémoire, de vacance et de rayonnement}

La conservation évolutive s'exprime mieux par l'échange que par l'immobilité scalaire. Un état complet répartit l'action et la quantité de mouvement entre mouvement visible, mémoire, déformation radiale, vacances, rayonnement et branche conjuguée de retour. Le bilan enrichi peut s'écrire

\[
\Delta J_{\mathrm{visible}}
+\Delta J_{\mathrm{memoria}}
+\Delta J_{\mathrm{radial}}
+\Delta J_{\mathrm{vacancias}}
+\Delta J_{\text{rayonnement}}
+\Delta J_{\mathrm{retorno}}
=0.
\]

Une force est \textbf{autoconservative} lorsque sa propre loi d'actualisation contient les canaux qui transportent, stockent, publient et réabsorbent l'action. La conservation relève du bilan généalogique complet de l'interaction. Chaque lecteur observe une partie de l'échange ; la généalogie recompose sa totalité.

Le système Terre–Lune offre une image physique directe. La rotation terrestre, l'orbite lunaire, les marées, la déformation et le rayonnement échangent du moment cinétique. La consommation nette du système complet se clôt dans le bilan généalogique de ces transferts. La grandeur appelée énergie fonctionne comme lecture dérivée de l'action et de la quantité de mouvement publiées dans chaque secteur. La conservation fondamentale réside dans la compatibilité entre tous les lecteurs de l'échange.

Cette formulation dissout l'ancienne frontière entre forces conservatives et forces dites non conservatives au moyen d'une catégorie plus profonde. Frottement, dissipation apparente et rayonnement occupent des canaux explicites de l'état élargi. Une perte dans la projection mécanique visible correspond à un transfert vers la mémoire, la déformation, la chaleur, les vacances ou le rayonnement. L'autoconservation suit la résidence de l'action à travers ces changements.

\subsection{Systèmes auto-inertiels : publication de l'accélération par orientation, courbure, feuillet et mémoire}

L'expression appropriée est \textbf{auto-inertie relationnelle}. Elle désigne la compatibilité entre la connexion propre de l'état total et la réponse qui apparaît lorsqu'on le projette dans la distribution globale. Les aspects endogène et exogène sont des composantes analytiques d'une seule relation dynamique.

Soit \(\Gamma\) une trajectoire horizontale ou géodésique de l'état complet. Sa loi propre satisfait

\[
\nabla_{\dot\Gamma}\dot\Gamma=0.
\]

En la projetant sur un lecteur \(S\), l'accélération observée est déterminée par la courbure de la projection :

\[
\nabla^S_{\dot\Gamma_S}\dot\Gamma_S
=
\mathcal K_{\pi_S}(\dot\Gamma,\dot\Gamma).
\]

L'accélération publiée exprime la relation entre mouvement total et géométrie du lecteur. Le principe de Mach complète ce mécanisme par la trace globale de frontière

\[
b:\mathcal X_\infty\longrightarrow B_\infty
\]

et la factorisation de la réponse inertielle locale

\[
I_v=\bar I_v\circ b.
\]

La distribution globale entre dans la réponse locale par une application typée. L'auto-inertie relationnelle réunit donc la connexion endogène du mouvement, la condition globale transmise par \(b\) et la projection concrète de l'observateur. Une écriture synthétique de cette coordination est

\[
\nabla^{\mathrm{auto}}
=
\mathcal M\bigl(
\nabla^{\mathrm{end}},
b[\mathcal X_\infty],
\pi_S
\bigr),
\]

où \(\mathcal M\) représente le couplage machien et conserve le typage de chaque composante.

Le principe d'Ambivalence ajoute la double orientation de la lecture. Toute publication directe conserve une branche conjuguée de retour ; émission et absorption forment les extrémités liées d'une même biographie ; l'orientation temporelle se lit depuis les deux sens du transport. La théorie de l'absorbeur trouve ici une formulation naturelle : source et récepteur clôturent conjointement le bilan généalogique de phase et d'action. Le présent torsionnel est l'interface où les deux orientations sont rendues compatibles.

TRIT agit à ce niveau comme cristal de temps : il organise la périodicité locale de l'actualisation et conserve, par la mémoire, la différence entre retours successifs. La branche directe et la branche conjuguée relient émission et absorption, et le couplage observateur–observable acquiert une résidence dynamique au sein du même état. Le temps publié est la lecture de cette orientation ; le temps généalogique est la séquence complète de phase, de retour et de mémoire.

Mach et Ambivalence remplissent des fonctions complémentaires. Mach relie la réponse locale à la distribution globale. Ambivalence conserve les directions directe et conjuguée de l'évolution. L'auto-inertie relationnelle articule les deux et remplace la scission entre systèmes inertiels et non inertiels par une loi unique de connexion, de projection et de mémoire.

\subsection{Dualité surfacique--volumique, transition hexade--octade et fonction du paramètre de Barbero--Immirzi}

Le passage de l'aire au volume constitue une transduction dimensionnelle centrale. L'aire publie l'action de frontière ; le volume publie la profondeur intérieure. La dérivation HMT du paramètre de Barbero–Immirzi organise la compatibilité entre les deux mesures par l'incidence hexade–octade et les canaux \(90/120\) :

\[
90=6\binom62,
\qquad
120=8\binom62.
\]

Les cardinaux \(6\) et \(8\) admettent une réalisation cubique par les faces et les sommets, tandis que l'hexade et l'octade de Steiner constituent des objets d'incidence distincts, reliés uniquement par le morphisme typé \(6\to8\). Les \(12\) arêtes du cube conservent leur propre fonction combinatoire et ne définissent pas à elles seules l'incidence de Steiner. La différence d'incidence exprime une incompatibilité dimensionnelle productive : l'information de surface nécessite une règle de transduction pour acquérir une résidence volumique.

Le même principe apparaît dans la retenue \(999+1=1000\). Les neuf atteignent le bord du registre décimal ; l'incrément publie simultanément la clôture du niveau précédent et l'origine \(0/1\) du suivant. Le rapport holographique entre \(729\), \(243\), \(27\) et \(1\) déploie cette transition aux échelles triadiques. La transduction de Barbero–Immirzi s'insère dans cette famille de changements de lecteur : frontière et intérieur conservent l'unité tout en redistribuant sa mesure.

Le défaut normalisé \(-1/12\) et le poids orienté \(12:-1\) remplissent des fonctions distinctes dans cette géométrie. Le premier quantifie une correction d'incidence ; le second enregistre une orientation signée dans la transduction. Leur coordination avec l'aire minimale, l'hamiltonien modulaire et la mémoire fournit le pont vers la dynamique dimensionnelle.

\subsection{Contraction hyperbolique intérieure, expansion extérieure et conservation holographique de la frontière}

La thèse peut se formuler à partir d'une idée unique : l'unité conserve son identité tout en changeant d'échelle, de lecteur, de courbure, de mémoire et de dimension. Cette permanence possède un caractère évolutif parce que l'unité se déploie, acquiert de nouvelles coordonnées et publie des propriétés différentes ; elle possède un caractère holographique parce que chaque publication finie conserve la règle de reconstruction de l'état dont elle procède. La clôture holographique de l'infini désigne précisément cette compatibilité entre identité et déploiement : chaque stade contient la loi qui permet d'approfondir vers l'intérieur et de prolonger vers l'extérieur la même généalogie.

La dénomination directrice est \textbf{principe de conservation évolutive de l'unité par transduction dimensionnelle et projection holographique}. L'expression situe chaque terme à son niveau correct. La conservation relève de l'identité généalogique ; l'évolution relève de l'actualisation ; la transduction relie des réalisations dimensionnelles ; la projection holographique conserve la correspondance reconstructive entre leurs lecteurs.

L'arithmétique élémentaire offre la première image de cette architecture. La complétion décimale

\[
1000=10^3=729+243+27+1=3^6+3^5+3^3+3^0
\]

expose quatre strates triadiques au sein d'une unité cubique décimale. En normalisant,

\[
1=\frac{729}{1000}+\frac{243}{1000}+\frac{27}{1000}+\frac{1}{1000},
\]

l'unité se reconnaît simultanément comme totalité et comme déploiement gradué. La lecture réciproque

\[
3^{-6},\quad 3^{-5},\quad 3^{-3},\quad 3^0
\]

décrit la profondeur intérieure de ces strates ; sa normalisation produit la partition correspondante. La relation \(999+1=1000\) exprime le mécanisme de transport : la dernière unité complète un registre et agit, en même temps, comme retenue inaugurant le suivant. Le zéro et le un forment ainsi une interface de continuité. Les intervalles \([0,1]\), \([0,10]\) et la prolongation vers \([0,\infty)\) sont des lecteurs scalaires d'une même loi de complétion.

La première scission \(1000=729+271\), suivie de \(271=243+27+1\), donne forme au rapport d'extrême et moyenne raison holographique du déploiement : un noyau triadique dominant conserve dans le reste la même règle de résolution. L'unité cubique décimale contient ainsi une profondeur triadique emboîtée. Le mouvement vers \(729\), \(243\), \(27\) et \(1\) lit la contraction intérieure ; la retenue vers le millier lit l'expansion extérieure.

Cette double orientation organise l'infini HMT. Vers l'intérieur, le raffinement croît hyperboliquement : chaque cellule admet une nouvelle profondeur sans perdre sa filiation. Vers l'extérieur, l'unité s'étend par des publications dimensionnelles successives. La clôture consiste en la commutativité des deux mouvements. Si \(U\) représente l'actualisation de l'état, \(R_d\) le lecteur dans la strate \(d\) et \(T_{d\to d'}\) la transduction entre strates, la conservation évolutive s'exprime par

\[
R_{d'}\circ U
=
T_{d\to d'}\circ R_d.
\]

Le même événement peut offrir une mesure angulaire, une densité d'action, une masse publiée ou une réponse gravitationnelle. Leur identité commune réside dans la généalogie qui fait commuter ces lectures.

\subsubsection{Monodromie nonadique et relèvement hélicoïdal avec avance de mémoire}

APP, TRIT et TPK fournissent le mécanisme formel de ce déploiement. APP construit le support et sépare les évaluations qui doivent être conservées. TRIT organise les régimes locaux et leur ambivalence. TPK transporte phase et mémoire sur le support enrichi. La dépendance causale est fixée par

\[
U_t\longrightarrow \Gamma_9\longrightarrow K_{\mathrm{ph}}.
\]

Le retour de phase et l'avance de mémoire forment des opérations coordonnées. Après neuf pas, la même classe angulaire peut réapparaître,

\[
g(K+9)=g(K),
\]

tandis que le compteur de mémoire progresse,

\[
q_{\mathrm{mem}}\Gamma_9=q_{\mathrm{mem}}+1,
\qquad
B_{K+9}\neq B_K.
\]

La figure géométrique précise est un relèvement hélicoïdal. Une orbite fermée dans la projection angulaire se relève en une trajectoire ouverte lorsque chaque retour de phase incrémente la mémoire. La circonférence conserve le cycle ; l'hélice conserve le cycle et enregistre en outre la profondeur atteinte. La monodromie coinductive génère ainsi une profondeur infinie par une succession de retours localement fermés et globalement ascendants.

Cette image dépasse la valeur simplement illustrative. La roue angulaire constitue une section de l'hélice ; l'hélice constitue la biographie complète de la roue lorsque l'état transporte de la mémoire. Chaque tour retrouve une position de phase et, simultanément, inaugure un nouveau feuillet. L'infinitude procède de la compatibilité entre retour et relèvement : réitérer le cycle produit de la profondeur sans dissoudre l'identité orbitale.

Le certificat nonadique donne un contenu fini et vérifiable à cette architecture. La chaîne

\[
w_6\longrightarrow w_{12}\longrightarrow w_{18}\longrightarrow
w_{24}\longrightarrow w_{30}\longrightarrow R_{36}\longrightarrow G_9
\]

ordonne le relèvement ; les \(396\) niveaux, les \(2376\) trits, les \(43\) tours et les \(387\) paires \((K,K+9)\) par canal quantifient le transit ; les cinq régions de \(\pi\) et la fibre ambiante de \(729\) éléments publient des lectures internes ultérieures. Cette séquence rend visible une monodromie productive : le retour retrouve la phase et la coinduction conserve la croissance de mémoire.

\subsubsection{Unité radionale et coordination des lecteurs circulaire, elliptique, projectif et hélicoïdal}

Le radion concentre dans une unité métrologique la structure que le radian publie angulairement. Sa définition naturelle est celle d'une \textbf{section orientée de phase, d'action et de mémoire}. Une unité de phase \(\theta=1\) transporte une unité d'action \(\sigma\hbar_{\mathrm{ret}}\), appartient à un tour complet \(T_+=2\pi_{\mathrm{HMT}}\) et conserve feuillet, orientation, retour nonadique et position dans le cylindre de raffinement.

Le radian exprime la coordonnée angulaire visible. Le radion conserve, avec cette coordonnée, l'action transportée, l'orientation, la retenue et la mémoire du tour. La métrologie angulaire acquiert ainsi une généalogie. Mesurer un arc signifie identifier quelle rotation a été publiée et quel état complet a réalisé cette rotation.

La relation entre phase et aire prend une forme particulièrement claire :

\[
\mathcal A(\theta)=\hbar_{\mathrm{ret}}\,\theta.
\]

Une unité radionale balaie une aire d'action \(\hbar_{\mathrm{ret}}\) ; un tour complet enferme \(h_{\mathrm{ret}}\). Le passage de la phase à l'aire est intégré à l'unité de lecture elle-même. La métrique angulaire, l'action et la mémoire cessent d'apparaître comme des quantités juxtaposées et deviennent les coordonnées d'une même section orientée.

Le radion admet plusieurs publications géométriques coordonnées. Le lecteur circulaire publie la phase ; le lecteur elliptique publie l'action et la forme ; le lecteur intérieur de Klein publie la profondeur ; le lecteur hélicoïdal nonadique publie le retour et la mémoire. L'unité reste généalogiquement identique à travers ces réalisations. C'est pourquoi le radion constitue le lien métrologique approprié entre HMT et la mécanique dimensionnelle : il offre une unité capable de traverser des changements de lecteur tout en conservant l'état que chaque mesure partielle présuppose.

\subsection{Polarité gravitationnelle à cinq composantes et transduction conjointe de l'action, de l'aire, du volume, du temps et de la masse}

La mécanique dimensionnelle permet de définir la dimension comme un régime stable de publication caractérisé par le rang minimal de mémoire indépendante requis par le relèvement. Traverser une dimension signifie appliquer un transducteur qui conserve la généalogie et change le type de lecture. La gravité occupe le niveau de connexion entre ces régimes.

Dans cette formulation, la \textbf{gravitation interdimensionnelle à cinq polarisations corrélatives}, en abrégé \textbf{structure pentapolaire de la gravitation}, est la connexion qui transporte l'unité généalogique et laisse à chaque interface une torsion minimale. Une écriture compacte de cette trace est

\[
\Theta_d^{\min}
=
\nabla_{d+1}\circ T_{d\to d+1}
-
T_{d\to d+1}\circ\nabla_d.
\]

L'expression mesure la différence entre transporter après avoir dérivé et dériver après avoir transporté. Cette différence est la mémoire géométrique minimale du franchissement dimensionnel. La composition d'interfaces produit une holonomie gravitationnelle

\[
\mathfrak G_{d\to D}
=
\operatorname{Hol}
\bigl(
\Theta_d^{\min},
\Theta_{d+1}^{\min},
\ldots,
\Theta_{D-1}^{\min}
\bigr).
\]

L'objet fondamental est cette connexion interdimensionnelle. Un mode local de type graviton occupe le rang dérivé de publication quantifiée de la torsion ; la gravitation conserve sa résidence primaire dans la compatibilité entre dimensions.

Le terme pentapolaire rassemble les cinq polarisations consubstantielles de l'état continu :

\[
\mathbf G_d^{\pm}
=
\bigl(
G_{\mathrm{sp},d}^{\pm},
G_{\mathrm{sol},d}^{\pm},
G_{\mathrm{gau},d}^{\pm},
G_{\mathrm{coh},d}^{\pm},
G_{\mathrm{det},d}^{\pm}
\bigr).
\]

Les cinq composantes émergent corrélativement du même état enrichi. La polarité \(\pm\) enregistre concentration directe et publication conjuguée. La gravité traverse les strates dimensionnelles en transportant simultanément ces cinq lectures ; la torsion minimale conserve la trace de chaque franchissement.

La géométrie acquiert ainsi une image unifiée. Le régime elliptique organise les clôtures ; le parabolique organise les seuils ; l'hyperbolique organise la profondeur. La transduction dimensionnelle relie ces régimes aux réponses radiales \(p\), aux holonomies de spin et aux publications de masse. La structure pentapolaire de la gravitation exprime la cohérence globale de ce lien.

La théorie M fournit des écrans physiques ultérieurs pour cette connexion : dimensions supplémentaires, cordes, branes et dualité T réalisent des modes concrets de transduction. HMT apporte la généalogie APP–TRIT–TPK, l'état enrichi, la mémoire et la lecture bidirectionnelle ; l'interface avec la théorie M publie ces structures dans des géométries physiques typées. La dualité T échange des lecteurs dimensionnels compatibles et conserve la classe généalogique de l'état.

\subsubsection{Compatibilité des réalisations classique, relativiste, quantique et HMT--MD}

La mécanique classique publie trajectoire et quantité de mouvement. La relativité publie structure causale, métrique et courbure. La mécanique quantique publie phase, holonomie, vacances et amplitudes. La mécanique dimensionnelle fournit les transducteurs qui relient ces publications, et HMT conserve la généalogie discrète qui les soutient.

Dans la lecture classique, une orbite décrit la projection visible du mouvement. Dans la lecture relativiste, cette orbite est intégrée à la géométrie du lecteur. Dans la lecture quantique, la phase et le spin enregistrent l'holonomie de la section. Dans la lecture HMT–MD, les trois expressions appartiennent à une même biographie relevée : le radion quantifie phase, action et mémoire ; \(p\) classe la réponse radiale ; \(\tau\) classe la courbure ; l'holonomie publie le spin ; le transducteur dimensionnel publie la masse ; la connexion pentapolaire publie la gravitation.

La dissipation classique apparente acquiert une résidence dans le bilan généalogique élargi. L'accélération relativiste s'exprime par connexion et projection. L'indétermination quantique s'organise par lecteurs, vacances et branches d'Ambivalence. Le principe de Mach relie la réponse locale à la distribution globale. Le Landauer nul garantit la réversibilité informationnelle de l'état complet. Tous ces éléments convergent vers une conservation plus profonde : l'identité de l'unité à travers les changements de publication.

\subsubsection{Théorème narratif de conservation généalogique et de prolongation holographique}

La clôture peut désormais se raconter comme une seule séquence. L'unité se décompose sans perdre son identité. La phase orientée acquiert une mesure dans le radion. Le retour nonadique transforme le cycle en hélice par la mémoire. Le caractère \(p\) transforme l'hélice en familles radiales et classe ses régimes. La persistance de ces trajectoires produit des sections ayant spin, charge, saveur, couleur et masse. Les vacances ouvrent des canaux ; la polarité électrique transporte la phase ; la réponse magnétique courbe transversalement ; les seuils publient le rayonnement. L'autoconservation clôt le bilan généalogique d'action et de quantité de mouvement. L'auto-inertie relationnelle relie connexion propre, projection et totalité machienne. L'Ambivalence conserve les directions directe et conjuguée. La transduction aire–volume conduit à Barbero–Immirzi et à la loi des masses. La torsion minimale enchaîne les dimensions et publie la structure pentapolaire de la gravitation.

La croissance vers l'intérieur et l'expansion vers l'extérieur sont des orientations conjuguées du même processus. Le raffinement hyperbolique génère la profondeur ; la transduction dimensionnelle génère l'extension. La retenue \(0/1\) relie les échelles ; la monodromie relie les tours ; la mémoire relie les états ; l'holonomie relie les orientations ; la gravitation relie les dimensions. Chaque lien conserve une unité qui évolue et chaque publication finie contient la règle de sa prolongation.

La thèse fondamentale est ainsi formulée :

\begin{quote}
\textbf{La clôture holographique de l'infini réalise le principe de conservation évolutive de l'unité par transduction dimensionnelle et projection holographique : l'unité maintient son identité généalogique en se déployant comme phase, action, courbure, mémoire, particule, champ et gravitation ; chaque lecture finie conserve la règle de reconstruction et de prolongation de l'état complet.}
\end{quote}

Cette formulation offre une image physique et mathématique conjointe. L'infini acquiert une loi de clôture locale ; l'unité acquiert de la profondeur ; la dimension acquiert de la mémoire ; la conservation acquiert un caractère relationnel ; la force acquiert l'autoconservation ; l'inertie acquiert une structure machienne ; la gravité acquiert une connexion pentapolaire ; et l'holographie acquiert le sens précis de reconstruction entre lecteurs compatibles.


\section{Définition opératoire du radion angulaire au moyen de l'état enrichi et de \(4\) lecteurs typés}
\subsection{Définition opératoire du radion angulaire et quadruple réalisation typée de l'état enrichi}
% Fragmento público generado desde propietarios íntegros.
% Residencia: 52.13.
% Fuente: ../incoming/radion_monografia_20260827/manuscrito/sections/13_sintesis.tex:1-80
\noindent La synthèse opératoire réunit les lecteurs de phase, d'action, de profondeur et de mémoire sans les identifier entre eux.\par\medskip

\subsubsection{Quadruple réalisation de l'état radional}

Soit \(x_\infty\in X_{\mathrm{enr}}\) la section coinductive. Le radion complet n'est pas
un point isolé, mais le paquet
\begin{equation}
\mathfrak R_\sigma(x_\infty)
=\bigl(
\theta=1,\sigma,
\pih,\alphah,A,C^*,\Delta_4,
\Phi_T,D_E,\epsone,\hret,
q_{\mathrm{mem}},\mathcal I_\infty
\bigr).
\label{eq:full-radion}
\end{equation}
Quatre lecteurs en extraient des objets différents :

\begin{figure}[H]
\centering
\begin{tikzpicture}[
  node distance=17mm and 18mm,
  state/.style={draw=RadNavy,very thick,rounded corners,fill=RadCream,
    minimum width=39mm,minimum height=13mm,align=center},
  readout/.style={draw=RadTeal,rounded corners,fill=RadMist,
    minimum width=39mm,minimum height=12mm,align=center},
  arr/.style={-{Latex},thick,RadCopper}]
\node[state] (x) {état enrichi\\\(x_\infty\)};
\node[readout,above left=of x] (circ) {cercle \(S^1\)\\phase};
\node[readout,above right=of x] (ell) {ellipse positive\\action et forme};
\node[readout,below left=of x] (klein) {interior Klein\\profundidad};
\node[readout,below right=of x] (helix) {hélice nonadique\\mémoire et retour};
\draw[arr] (x)--(circ) node[midway,below left]{\small \(P_{\mathrm{ph}}\)};
\draw[arr] (x)--(ell) node[midway,below right]{\small \(P_{\mathrm{act}}\)};
\draw[arr] (x)--(klein) node[midway,above left]{\small \(P_K\)};
\draw[arr] (x)--(helix) node[midway,above right]{\small \(P_{\mathrm{mem}}\)};
\end{tikzpicture}
\caption{Quatre publications typées du même antécédent. L'unité de
l'état n'élimine pas la différence des codomaines.}
\label{fig:four-readers}
\end{figure}

Le cercle conserve la classe de phase et oublie le relèvement. L'ellipse conserve
l'action, l'anisotropie et l'orientation symplectique. Klein dote l'intérieur d'une
distance projective. L'hélice conserve le nombre de retours et le raffinement.
La mécanique dimensionnelle ne consiste pas à les mélanger, mais à construire les
transducteurs qui les rendent compatibles.

\subsubsection{Interprétation opératoire du cocycle \(\varepsilon_R\)}

\(\varepsilon_R\) n'est pas une constante universelle indépendante de tout contexte. C'est la
coordonnée d'un cocycle de transition associé à la carte du radion. Son
domaine contient deux sections du même état :
\[
S_E=1+D_E,
\qquad
S_T=1+\Phi_T.
\]
Sa règle est
\begin{equation}
\epsone=S_T-S_E=\Phi_T-D_E.
\label{eq:epsilon-final-meaning}
\end{equation}
Son codomaine est la droite adimensionnelle d'unité angulaire ; la carte
\(360/T_+\) la publie en degrés.

\begin{exactbox}
\textbf{En langage simple :} la section directe et la section torsionnelle arrivent
à la même cellule d'unité par deux chemins internes différents. Leurs parties
entières coïncident ; leurs restes, non. \(\epsR\) est la différence orientée de ces
restes, calculée avec la retenue d'APP et portée à la limite par les cylindres
TPK. Elle n'est pas introduite pour que le radian se ferme : elle apparaît avant de consulter le
radian et, lorsqu'elle est publiée, elle le ferme.
\end{exactbox}

Cette définition explique pourquoi le nombre est petit. Non parce qu'on a
cherché une petite correction, mais parce que \(S_E\) et \(S_T\) partagent un
préfixe profond. La petitesse mesure la proximité structurelle de deux
publications ; le signe mesure leur orientation ; les chiffres enregistrent la retenue.


% Fragmento público generado desde propietarios íntegros.
% Residencia: 52.13.1.
% Fuente: ../incoming/radion_monografia_20260827/manuscrito/sections/13_sintesis.tex:81-134
\subsubsection{Définition opératoire du radion}

\begin{definition}[Radion HMT--MD]
Le radion est la section orientée de phase, d'action et de mémoire qui satisfait
simultanément :
\begin{enumerate}
\item parcourt une unité de phase \(\theta=1\) ;
\item balaie une action signée \(\sigma\hret\) ;
\item appartient à un tour \(T_+=2\pih\) généré coinductivement ;
\item conserve le feuillet \(\sigma\), le retour nonadique et le cylindre de
raffinement ;
\item satisfait le raccord exact
\(1=S_E-\Phi_T+\epsone\).
\end{enumerate}
\end{definition}

Le radian conventionnel est l'image du radion par le foncteur d'oubli
\[
\mathfrak R_\sigma
\longmapsto
\theta=1\in\R/2\pi\Z.
\]
Cet oubli est utile et légitime pour la trigonométrie euclidienne. Il devient
insuffisant lorsque la physique doit distinguer action, signe, retour,
champ et histoire.

\paragraph{Généalogie de l'unité angulaire naturelle}

Le caractère naturel du radian est généralement justifié par le fait que la dérivée de
\(e^{i\theta}\), \(\sin\theta\) et \(\cos\theta\) adopte sa forme la plus simple
lorsque \(\theta\) est mesuré comme quotient arc/rayon. HMT conserve cet avantage,
mais l'explique depuis une couche antérieure :
\[
\theta\quad\leftrightarrow\quad
\frac{\mathcal A(\theta)}{\hret}.
\]
La phase est naturelle parce qu'elle mesure l'action en unités de \(\hret\) ; le tour est
naturel parce qu'il enferme \(\hfull\). Le quotient arc/rayon est la publication
euclidienne d'une unité qui possède déjà une structure symplectique et une mémoire.

\paragraph{Compatibilité entre ellipse d'action et croissance dimensionnelle}

Le cercle est la forme qui résulte de l'oubli de l'anisotropie :
\(C^*=0\Rightarrow\eta=0\Rightarrow a=b\). L'ellipse est plus informative parce qu'elle
conserve deux valeurs propres, deux axes et une orientation d'échange à produit
fixe. En ce sens, l'ellipse est une prolongation du cercle, non un cercle
déformé par accident.

La profondeur de Klein et le relèvement hélicoïdal ajoutent deux types d'information
que le contour plan ne contient pas : distance intérieure et mémoire de retour.
Telle est l'image ontologique de la mécanique dimensionnelle qui émerge du radion :
une dimension nouvelle apparaît lorsque l'on exige de conserver une donnée indépendante
que la projection précédente devait oublier.


\subsection{Propriétés structurelles de l'unité radionale : généalogie, charnière critique, densité, retenue, action, compatibilité bord--intérieur, forme, polarité et modularité}
% Fragmento público generado desde propietarios íntegros.
% Residencia: 52.13.2.
% Fuente: ../incoming/radion_monografia_20260827/manuscrito/sections/13_sintesis.tex:135-184
\subsubsection{Système de théorèmes composants du radion}

\begin{theorem}[Généalogie]
L'objet \(\mathfrak R_\sigma\) descend de
\(\APP\to\TRIT\to\TPK\to X_{\mathrm{enr}}\) et ne reçoit pas la cible métrologique
comme sélecteur.
\end{theorem}

\begin{theorem}[Charnière]
Dans la carte APP de cent marques, la couture critique et la deuxième phase
dodécaphasique satisfont \(j_{100}(1/2)=\theta_2=50^\circ\), avec \(m=2\) unique.
\end{theorem}

\begin{theorem}[Densité]
Les deux feuillets du secteur actif \(N_6=90\), normalisés par \(3^6=729\),
produisent \(\rho_{\mathrm{tw}}=20/81\).
\end{theorem}

\begin{theorem}[Retenue]
La soustraction APP des sections \(S_T\) et \(S_E\) converge vers
\(\epsone\) avec un taux d'intervalle \(3^{-54}\) par retour et ferme exactement
l'unité.
\end{theorem}

\begin{theorem}[Action]
L'ellipse positive satisfait \(\mathcal A(\theta)=\hret\theta\) ; un radion
balaie \(\hret\) et un tour enferme \(\hfull\).
\end{theorem}

\begin{theorem}[Bord--intérieur]
La même frontière possède une action finie \(\hfull\), tandis que l'intérieur de
Klein a une aire hyperbolique infinie.
\end{theorem}

\begin{theorem}[Forme]
Le quotient \(e^{-\eta}\) est conservé à travers les demi-axes spectraux,
les demi-axes d'action, les incertitudes, l'impédance et le module du tore.
\end{theorem}

\begin{theorem}[Polarité]
L'inversion \(C^*\mapsto-C^*\) échange les lois constitutives, inverse
l'impédance et conserve la vitesse normalisée.
\end{theorem}

\begin{theorem}[Modularité]
L'inversion bidirectionnelle agit comme \(S:\tau\mapsto-1/\tau\) et conserve
\(j(\tau)\), tandis que l'étiquetage focal conserve l'orientation que
\(j\) oublie.
\end{theorem}

% Fuente: ../incoming/radion_monografia_20260827/manuscrito/sections/A_demostraciones.tex:1-179
\subsubsection{Démonstrations algébriques et géométriques réunies}

\paragraph{Démonstration linéaire de la clôture unitaire}

Pour faciliter un audit indépendant, nous réunissons l'argument sans
interruption narrative.

\begin{enumerate}[label=\textbf{Étape \arabic*.},leftmargin=2.2cm]
\item La connexion nonadique génère de manière corrélée
\(\pih,e_{\HMT},\varphi_{\HMT},\alphah\).
\item Le tour est \(T_+=2\pih\).
\item La roue fixe \(\theta_2=50^\circ\), et la carte parangulaire fixe
\(A=1000\alphah\).
\item La section directe est
\[
S_E=T_+\left(\frac{50}{360}+\frac{1000\alphah}{360}\right)
=2\pih\left(\frac5{36}+\frac{25}{9}\alphah\right).
\]
\item Les intervalles certifiés donnent \(1<S_E<2\) ; donc
\(D_E=S_E-1\).
\item L'incidence produit \(N_6=90\) ; deux feuillets et l'espace ambiant \(729\)
produisent \(\rho=2N_6/729=20/81\).
\item La torsion globale est
\[
\Delta_4=\frac{\pih-e_{\HMT}\log\pih}{270}
\left(1+\frac{\pih}{729}\right).
\]
\item La section torsionnelle est \(S_T=1+(20/81)\Delta_4\).
\item La soustraction APP définit
\[
\epsone=S_T-S_E=\frac{20}{81}\Delta_4-D_E.
\]
\item Par substitution,
\[
S_E-\frac{20}{81}\Delta_4+\epsone=1.
\]
\item Multiplier par \(360/T_+=180/\pih\) produit
\[
\frac{180^\circ}{\pih}
=50^\circ+A^\circ-\frac{20}{81}\Delta_4^\circ+\epsR^\circ.
\]
\end{enumerate}

Aucune étape ne reçoit le membre de gauche de la dernière égalité. Le résultat
métrologique n'est évalué qu'à la fin comme reconnaissance.

\paragraph{Unicité de la soustraction avec retenue}

Soit \(B=1000\). Pour chaque position, tout entier \(u_i\) admet une division
euclidienne unique
\[
u_i=Bc_i+r_i,\qquad 0\le r_i<B.
\]
Par conséquent,
\[
r_i=u_i\bmod B,\qquad c_i=(u_i-r_i)/B
\]
sont uniques. En parcourant les positions de l'extrémité distante vers
l'avant, la retenue de sortie d'une position fixe l'entrée de la suivante.
Par récurrence, tout le mot de restes et de retenues est unique.

En pseudocode :
\begin{verbatim}
acarreo = remote_acarreo
for i from last_block downto first_block:
    u = torsion[i] - direct[i] + acarreo
    remainder[i] = u mod 1000
    acarreo = (u - remainder[i]) / 1000
return remainder, acarreo_word
\end{verbatim}
La procédure ne consulte pas la valeur attendue de la différence.

\paragraph{Convergence de la fonctionnelle selon la profondeur}

Soit
\[
F(p,e)=1+\frac{20}{81}
\frac{p-e\log p}{270}\left(1+\frac{p}{729}\right)
-2p\left(\frac5{36}+\frac{25}{9}\alphah\right).
\]
Dans le rectangle compact \(3\le p\le4\), \(2\le e\le3\), \(F\) est
continue et continûment différentiable. Si \(I_N^\pi\searrow\{\pih\}\) et
\(I_N^e\searrow\{e_{\HMT}\}\), l'image par intervalles
\[
F(I_N^\pi,I_N^e)
\]
forme une famille emboîtée dont la largeur tend vers zéro. L'intersection est un
unique réel, \(\epsone\). La monotonie des dérivées permet d'évaluer les
extrémités sans échantillonnage intérieur.

Le taux \(3^{-54}\) par retour dérive du raffinement de six trits sur
neuf niveaux. À la profondeur de \(45\) triades, la borne de largeur
\(<4.81\times10^{-130}\) certifie plus de 129 chiffres de la valeur.

\paragraph{Dérivation de la métrique et de l'aire hyperbolique de Klein}

Dans le disque unité, la forme matricielle de la métrique de Klein est
\[
g(u,v)=\frac1{1-r^2}I
+\frac1{(1-r^2)^2}
\begin{pmatrix}u\\v\end{pmatrix}
\begin{pmatrix}u&v\end{pmatrix}.
\]
Le lemme du déterminant pour une perturbation de rang un donne
\[
\det g=(1-r^2)^{-3}.
\]
D'où
\[
\sqrt{\det g}\,\diff u\diff v
=(1-r^2)^{-3/2}\diff u\diff v.
\]
L'intégration en coordonnées polaires produit \eqref{eq:klein-area-R}. La limite infinie est
due à la singularité métrique à la frontière ; la frontière conserve une
aire symplectique finie parce que cette mesure utilise \(\diff Q\wedge\diff P\),
non \(\sqrt{\det g}\).

\paragraph{Distance focale dans le disque de Klein}

Sur un diamètre, la distance du centre à \(r\) est
\[
d_K(0,r)=\operatorname{artanh}r.
\]
Pour \(r=e_f\), \eqref{eq:excentricidad} implique
\[
1-e_f^2=e^{-2\eta}.
\]
Si \(u_f=\operatorname{artanh}e_f\), alors
\[
\operatorname{sech}u_f=\sqrt{1-e_f^2}=e^{-\eta},
\]
et donc \(\eta=\log\cosh u_f\), comme dans
\eqref{eq:focal-hyperbolic-chain}.

\paragraph{Transformation symplectique de compression}

Définissons
\begin{equation}
S_\eta=
\begin{pmatrix}e^{\eta/2}&0\\0&e^{-\eta/2}\end{pmatrix},
\qquad \det S_\eta=1.
\label{eq:compresión simpléctica}
\end{equation}
La structure complexe transportée est
\begin{equation}
J_\eta=S_\eta
\begin{pmatrix}0&-1\\1&0\end{pmatrix}
S_\eta^{-1}
=\begin{pmatrix}0&-e^\eta\\e^{-\eta}&0\end{pmatrix},
\qquad J_\eta^2=-I.
\label{eq:Jeta}
\end{equation}
La courbe
\[
\gamma(\theta)=\sqrt{2\hret}\,
S_\eta(\cos\theta,\sin\theta)
=e^{\theta J_\eta}\gamma(0)
\]
est exactement l'ellipse d'action. La compression symplectique hyperbolique conserve l'aire et
transporte la rotation complexe au lieu de la détruire.

Pour l'orientation \(\sigma\),
\[
\gamma_\sigma(\theta)
=(a_S\cos\theta,\sigma b_S\sin\theta),
\qquad
\mathcal A_\sigma(\theta)=\sigma\hret\theta.
\]
Les deux feuillets partagent les foyers et le module d'action ; ils diffèrent par l'action signée.

\paragraph{Vérification diophantienne du poids dodécaphasique de Barbero--Immirzi}

L'équation \(4a-3b=45\) a pour solutions entières
\[
(a,b)=(12+3t,1+4t),\qquad t\in\Z.
\]
La condition \(|b|=1\) laisse \(b=1\), car \(b=-1\) n'appartient pas à la classe
\(1\pmod4\). D'où \(a=12\). Cette vérification reproduit le couple obtenu par la chaîne fondamentale
dodécaphasique : douze secteurs et une unique couture orientée. Les cardinaux
\(90/120\), les multiplicités de la chaîne et le lecteur de retour
conservent des fonctions distinctes ; le coefficient de pôle et la minimalité
vérifient arithmétiquement leur composition.

\subsection{Sensibilité de la clôture unitaire aux variations contrôlées de la multiplicité des feuillets, de la phase dodécaphasique, du canal \(e\), de la torsion globale, de l'opérateur harmonique et de la queue sexadique}
% Fragmento público generado desde propietarios íntegros.
% Residencia: 52.13.3.
% Fuente: ../incoming/radion_monografia_20260827/manuscrito/sections/13_sintesis.tex:185-247
\subsubsection{Contrôles négatifs du système radional}

La construction est réfutable en des points concrets :
\begin{enumerate}
\item Si \(S_E\notin(1,2)\), le quotient APP n'est pas un et la clôture change de
cellule.
\item Si l'incidence active ne produit pas \(N_6=90\), ou s'il n'y a pas deux feuillets,
\(20/81\) ne s'ensuit plus.
\item Si le vérificateur reçoit \(180/\pi\) ou \(\epsR\) avant de construire les
opérandes, la génération est circulaire.
\item Si la soustraction par triades ne stabilise pas des préfixes compatibles, la
section coinductive proposée n'existe pas.
\item Si \(|C^*|\ge A\), l'ellipse positive cesse d'exister.
\item Si \(a_Sb_S\ne2\hret\), le théorème d'aire échoue.
\item Si \(C^*=0\), on doit avoir \(d=0\), \(\eta=0\), \(\tau=i\) et
\(j=1728\).
\item Si l'aire de Klein du sous-disque ne suit pas \eqref{eq:klein-area-R}, la
thèse de profondeur infinie échoue.
\item Si \(F_{\mathrm{spec}}=\epsR^\circ\) est affirmé sans contre-terme indépendant, la
différence numérique \eqref{eq:chiR} le réfute.
\item La production de \(12:-1\) exige de conserver, avec les cardinaux
\(90/120\), les douze secteurs dodécaphasiques, l'unique couture orientée et
l'action du lecteur sur ces générateurs. Changer une multiplicité
doit modifier l'image \(12S_{90}-S_{120}\). Pour la chaîne fondamentale
sans altération, le coefficient de \((1-q)^{-1}\) doit donner \(1/24\).
\item Si la fonction \(j\) est utilisée pour sélectionner le signe de charge, on perd
l'invariance \eqref{eq:j-invariance} et l'on commet une erreur de type.
\end{enumerate}

\subsubsection{Réalisation physique conjointe du radion}

L'image finale n'est pas un cercle corrigé. C'est une architecture de
publications :
\[
\begin{aligned}
&\text{APP:}&&\text{unité, reste, retenue et fibre};\\
&\text{TRIT:}&&\text{régime, orientation et feuillets};\\
&\text{TPK:}&&\text{phase, retour, dodécaphase et mémoire};\\
&\text{ellipse :}&&\text{action et anisotropie};\\
&\text{Klein:}&&\text{profondeur intérieure};\\
&\text{hélice :}&&\text{biographie du retour};\\
&\text{MD:}&&\text{horloge, énergie, charge, champ et causalité}.
\end{aligned}
\]

L'équation du radion réunit la face visible :
\BigRadionEquation
L'ellipse réunit la face d'action :
\[
\mathcal A(1)=\hret,\qquad \mathcal A(2\pi)=\hfull.
\]
La métrique de Klein réunit la profondeur :
\[
K=-1,\qquad \operatorname{Area}_K=\infty.
\]
L'holonomie réunit la mémoire :
\[
g(k+9)=g(k),\qquad B_{k+9}\ne B_k.
\]

\begin{thesisbox}
Le cercle est l'ombre angulaire ; l'ellipse est le corps d'action ; Klein mesure
l'intérieur ; l'hélice conserve la biographie. Le radion est l'unité qui
permet de lire les quatre affirmations sur un unique état sans réduire l'une à
l'autre.
\end{thesisbox}
% Fuente: ../incoming/radion_monografia_20260827/manuscrito/sections/B_valores_y_verificacion.tex:1-105
\subsubsection{Valeurs canoniques, stabilité et vérification reproductible}

\paragraph{Valeurs des invariants internes générés}

\begin{longtable}{p{.27\textwidth} p{.61\textwidth}}
\toprule
Objet & Valeur utilisée dans le profil coinductif\\
\midrule
\endfirsthead
\toprule
Objet & Valeur utilisée dans le profil coinductif\\
\midrule
\endhead
\(\pih\) & \(\begin{aligned}3.1415926535897932384626433832795028841\\[-2pt]971693993751058209749445923\ldots\end{aligned}\)\\
\(e_{\HMT}\) & évaluation coinductive du lecteur d'Euler, sans table externe\\
\(\alphah\) & \(\begin{aligned}0.0072973525692838009972851054723806629\\[-2pt]995641725396427091854046825\ldots\end{aligned}\)\\
\(A\) & \(\begin{aligned}7.2973525692838009972851054723806629995\\[-2pt]641725396427\ldots{}^\circ\end{aligned}\)\\
\(C^*\) & \(2.1237383389686457242619513803933607937\ldots{}^\circ\)\\
\(\eta\) & \(0.299689683921630799684926562863927\ldots\)\\
\(\hret\) & \(1.05457181691503906113369058472430\ldots\times10^{-34}U_S\)\\
\bottomrule
\end{longtable}

\paragraph{Convergence quantitative selon la profondeur}

\begin{center}
\small
\begin{tabular}{r l l}
\toprule
Profondeur & \(\epsR^\circ\) & lecture\\
\midrule
12 & \(5.1681282186551375597\times10^{-8}\) & préasymptotique\\
18 & \(5.1383017764316644535\times10^{-8}\) & 7 cifras estables\\
24 & \(5.1383016758575394560\times10^{-8}\) & 14 cifras estables\\
30 & \(5.1383016758575282072\times10^{-8}\) & 20 cifras estables\\
36 & \(5.138301675857528207168517\times10^{-8}\) & perfil largo\\
45 & \(5.13830167585752820716851646226459\times10^{-8}\) & largeur \(<4.81\times10^{-130}\)\\
\bottomrule
\end{tabular}
\end{center}

Les retours de profondeur \(9,18,27,36,45\) conservent la phase \(8\) et
augmentent la mémoire \(0,1,2,3,4\). Le rapport entre largeurs consécutives tend
vers
\[
3^{-54}=1.7196982334549856127610399316004871\ldots\times10^{-26}.
\]

\paragraph{Analyse de sensibilité structurelle par suppression contrôlée}

\begin{longtable}{p{.46\textwidth} p{.34\textwidth}}
\toprule
Opération & Sortie ou différence unitaire\\
\midrule
\endfirsthead
\toprule
Opération & Sortie ou différence unitaire\\
\midrule
\endhead
Opérateur canonique & \(8.96802822044562981999\times10^{-10}\)\\
Un feuillet, \(\rho=90/729\) & \(-1.3727056615960678\times10^{-5}\)\\
Phase \(\theta_1=20^\circ\) & \(+5.2359877649510169\times10^{-1}\)\\
Phase \(\theta_3=80^\circ\) & \(-5.2359877470149605\times10^{-1}\)\\
Sans canal \(e\) & \(1.8065350748904653\times10^{-3}\)\\
Torsion historique \(\Delta_4(A,C^*)\) & \(1.0517084608768816\times10^{-9}\)\\
Opérateur harmonique \(M_4\) moins canonique & \(1.0525500222895070\times10^{-17}\)\\
Queue sexadique moins canonique & \(2.0283936357433839\times10^{-12}\)\\
\bottomrule
\end{longtable}

La séparation des échelles montre une sensibilité structurelle : la clôture dépend
de la phase, des deux feuillets, des canaux \(\pi/e\), de la torsion globale et de la retenue.
Aucune ablation ne reproduit la valeur.

\paragraph{Certificats informatiques reproductibles}

L'ouvrage utilise les propriétaires exécutables suivants :
\begin{enumerate}
\item \sourcepath{output/RADION_TWOWAY_NONADICO_20260827/verificar_emergencia_radion_por_profundidad.py};
\item \sourcepath{output/RADION_TWOWAY_NONADICO_20260827/verificar_radion_helice_critica_tpk.py};
\item \sourcepath{output/RADION_TWOWAY_NONADICO_20260827/verificar_operador_tpk_target_free_profundidad.py};
\item le vérificateur consolidé de cette monographie, situé dans
\sourcepath{certificados/verificar_monografia_radion.py}.
\end{enumerate}
Les sorties attendues comprennent
\[
\texttt{PASS\_EMERGENCIA\_RADION\_POR\_PROFUNDIDAD},
\]
\[
\texttt{PASS\_RADION\_HELICE\_CRITICA\_TPK},
\qquad
\texttt{PASS\_MONOGRAFIA\_RADION\_HMT\_MD}.
\]

\paragraph{Distinction entre précision informatique et exactitude mathématique}

Qu'un programme affiche zéro à 120 ou 200 chiffres ne transforme pas à lui seul une
identité résiduelle en prédiction. L'exactitude mathématique procède de
\eqref{eq:cierre-unitario} ; la précision arbitraire permet d'évaluer ses
opérandes et de vérifier la convergence. La non-circularité procède de l'ordre
de construction et de l'interdiction des cibles, non du nombre de chiffres.

Le calcul Decimal peut laisser des résidus d'ordre \(10^{-218}\) lorsqu'il combine
des branches évaluées avec des précisions différentes. Ce nombre est un arrondi de
l'implémentation ; non une correction physique supplémentaire.
% Fuente: ../incoming/radion_monografia_20260827/manuscrito/sections/D_historia_fuentes.tex:1-169
\subsubsection{Généalogie conceptuelle et sources primaires}

\paragraph{Quantum d'action de Planck}

Planck introduit \(h\) dans la loi de rayonnement comme constante reliant énergie
et fréquence. La donnée historique décisive pour cet ouvrage est dimensionnelle :
\(h\) est une action. Le radion ne modifie pas ce rôle ; il explique pourquoi l'action
d'un tour et l'action par unité de phase sont liées par \(2\pi\).

\paragraph{Quantification angulaire de Bohr et de Sommerfeld}

Bohr utilise \(h/2\pi\) dans la quantification du moment cinétique. Sommerfeld étend
la quantification aux intégrales d'action et aux orbites képlériennes. Il n'est donc pas
historiquement exact de dire que Dirac a inventé la constante réduite. Sa
contribution appartient à une généalogie dans laquelle l'action angulaire était déjà
présente.

La nouveauté HMT--MD est différente : elle ne prend pas \(h/2\pi\) comme un quotient externe,
mais construit une ellipse dont le tour a pour aire \(h\) et dont le secteur d'un
radian a pour aire \(\hbar\).

\paragraph{Formalisme matriciel de Born et de Jordan et oscillateur quantique}

Born et Jordan placent le principe variationnel, les phases et l'oscillateur
harmonique au cœur de la mécanique matricielle ; le champ électromagnétique se
décompose en une collection d'oscillateurs. La réalisation LC du radion
prolonge cette intuition avec une donnée supplémentaire : le quotient de forme
\(e^{-\eta}\) détermine l'impédance relative tandis que produit, fréquence
et action restent fixes.

\paragraph{Fréquence angulaire et action dans la formulation de Dirac}

En 1925, Dirac écrit les fréquences en radians par unité de temps et utilise
des relations de la forme
\[
xy-yx=\frac{ih}{2\pi}[x,y],
\qquad
\Delta E=\frac{h}{2\pi}\omega.
\]
L'équation \(E=\hret\omega\) relie explicitement action réduite et
fréquence angulaire. Dans le radion, cette relation reçoit une géométrie : \(\hret\)
est l'aire balayée par une unité de phase.

\paragraph{Relations d'indétermination de Heisenberg, Robertson et Schrödinger}

Heisenberg formule la limitation conjointe des variables canoniquement
conjuguées. Robertson généralise l'inégalité ; Schrödinger intègre la
covariance. La matrice \eqref{eq:covariance} relève plus directement de cette
prolongation Robertson--Schrödinger que de l'article de Heisenberg pris isolément.

Aucun de ces auteurs ne formule l'ellipse généalogique complète de cet ouvrage.
L'affirmation historiquement défendable est qu'ils fournissent le langage dans
lequel l'ellipse d'action peut être reconnue comme un contour minimal ; la
construction du couple \((A,C^*)\), les foyers, Klein, l'hélice et la retenue constituent la
contribution HMT--MD.

\paragraph{Fonctionnelle de phase de Feynman}

La formulation spatio-temporelle attribue à chaque chemin une phase
\(e^{iS/\hbar}\). L'unité de radion rend la correspondance littérale :
incrémenter l'action de \(\hbar\) incrémente la phase d'un radian. La somme des
chemins reconnaît, dans un langage intégral, le même principe que l'ellipse
exprime comme aire.

\paragraph{Propagations avancée et retardée chez Wheeler--Feynman}

La théorie de l'absorbeur construit une interaction classique temporellement
symétrique au moyen de combinaisons de champs retardés et avancés. L'identité
HMT \(CU(t)C^{-1}=U(-t)\) et l'inversion des chemins fournissent un
antécédent opératoire. La promotion physique exige que l'application entre chemins
et propagateurs préserve action, orientation, frontière et forme symplectique.

La théorie de l'absorbeur n'est pas la même affirmation que l'interprétation de
Feynman et de Stückelberg des antiparticules comme propagation temporellement
inversée. Les deux peuvent reconnaître la bidirectionnalité ; leurs objets physiques et
leurs conditions de frontière sont différents.

\paragraph{Dualité de Hodge et séparation d'avec la cohomologie de Hochschild}

L'opération pertinente dans ce manuscrit est l'étoile de Hodge
\(\star^2=-I\) sur les deux-formes en signature lorentzienne. On n'a pas localisé dans
le corpus actif du radion de propriétaire d'homologie de Hochschild qui
agisse sur l'ellipse, la retenue ou le secteur \(90/120\). Introduire Hochschild
par ressemblance nominale dégraderait le typage. Si une prolongation future
construit un cocycle de déformation avec ses propres domaine et codomaine, il devra
entrer comme résultat nouveau, non comme rétroattribution.

\paragraph{Bibliographie primaire sélectionnée}

\begin{enumerate}
\renewcommand{\labelenumi}{[\arabic{enumi}]}
\item\hypertarget{radionbib:Planck1901}{}
M.~Planck,
\emph{Ueber das Gesetz der Energieverteilung im Normalspectrum},
Annalen der Physik 309 (1901), 553--563.
\url{https://doi.org/10.1002/andp.19013090310}

\item\hypertarget{radionbib:Bohr1913}{}
N.~Bohr,
\emph{On the Constitution of Atoms and Molecules},
Philosophical Magazine 26 (1913), 1--25.
\url{https://doi.org/10.1080/14786441308634955}

\item\hypertarget{radionbib:Sommerfeld1916}{}
A.~Sommerfeld,
\emph{Zur Quantentheorie der Spektrallinien, I},
Annalen der Physik 356 (1916), 1--94.
\url{https://doi.org/10.1002/andp.19163561702}

\item\hypertarget{radionbib:BornJordan1925}{}
M.~Born et P.~Jordan,
\emph{Zur Quantenmechanik},
Zeitschrift für Physik 34 (1925), 858--888.
\url{https://doi.org/10.1007/BF01328531}

\item\hypertarget{radionbib:Dirac1925}{}
P.~A.~M.~Dirac,
\emph{The Fundamental Equations of Quantum Mechanics},
Proceedings of the Royal Society A 109 (1925), 642--653.
\url{https://doi.org/10.1098/rspa.1925.0150}

\item\hypertarget{radionbib:Heisenberg1927}{}
W.~Heisenberg,
\emph{Über den anschaulichen Inhalt der quantentheoretischen Kinematik und Mechanik},
Zeitschrift für Physik 43 (1927), 172--198.
\url{https://doi.org/10.1007/BF01397280}

\item\hypertarget{radionbib:Dirac1928}{}
P.~A.~M.~Dirac,
\emph{The Quantum Theory of the Electron},
Proceedings of the Royal Society A 117 (1928), 610--624.
\url{https://doi.org/10.1098/rspa.1928.0023}

\item\hypertarget{radionbib:Robertson1929}{}
H.~P.~Robertson,
\emph{The Uncertainty Principle},
Physical Review 34 (1929), 163--164.
\url{https://doi.org/10.1103/PhysRev.34.163}

\item\hypertarget{radionbib:Schrodinger1930}{}
E.~Schrödinger,
\emph{Zum Heisenbergschen Unschärfeprinzip},
Sitzungsberichte der Preussischen Akademie der Wissenschaften (1930), 296--303.

\item\hypertarget{radionbib:WheelerFeynman1945}{}
J.~A.~Wheeler et R.~P.~Feynman,
\emph{Interaction with the Absorber as the Mechanism of Radiation},
Reviews of Modern Physics 17 (1945), 157--181.
\url{https://doi.org/10.1103/RevModPhys.17.157}

\item\hypertarget{radionbib:Feynman1948}{}
R.~P.~Feynman,
\emph{Space-Time Approach to Non-Relativistic Quantum Mechanics},
Reviews of Modern Physics 20 (1948), 367--387.
\url{https://doi.org/10.1103/RevModPhys.20.367}

\item\hypertarget{radionbib:Feynman1949}{}
R.~P.~Feynman,
\emph{The Theory of Positrons},
Physical Review 76 (1949), 749--759.
\url{https://doi.org/10.1103/PhysRev.76.749}
\end{enumerate}

\begin{narrativebox}
Les fondateurs contiennent les éléments : action, fréquence angulaire, variables
conjuguées, oscillateur, incertitude, phase des chemins et réversion temporelle.
La contribution du radion consiste à les réunir dans un objet généré : une ellipse
d'aire \(h\), un secteur d'aire \(\hbar\), un intérieur de profondeur infinie,
une orientation de feuillet et une retenue exacte de compatibilité.
\end{narrativebox}
% Fuente: ../incoming/radion_monografia_20260827/manuscrito/sections/E_conclusion.tex:1-107
\subsubsection{Synthèse finale : unité angulaire avec généalogie de phase, d'action et de mémoire}

Le radian ne disparaît pas. Il devient entièrement visible.

Dans son usage ordinaire, un radian est un rapport qui n'a pas besoin de se souvenir de la façon dont il a été
produit. L'arc et le rayon s'annulent dimensionnellement ; l'unité semble
nue. Cette nudité est une vertu pour la trigonométrie, mais c'est un oubli
pour une théorie qui entend expliquer la structure du continu et l'origine
de l'action. Le radion récupère l'oublié sans modifier la valeur publiée.

La généalogie commence dans APP, où une sortie conserve la valeur et aussi
le feuillet, le quotient, le reste, la retenue et la frontière. Le TRIT introduit régime et
orientation. Le TPK fait revenir la phase sans réinitialiser l'état. La connexion
nonadique construit la section coinductive de \(\pi\), et la dodécaphase publie
la charnière de \(50^\circ\). Le couple angulaire apporte un mode commun \(A\) et une
ouverture orientée \(C^*\). L'incidence de six phases actives produit
\(90\), ses deux feuillets sur l'espace ambiant \(729\) produisent \(20/81\), et la
torsion globale apporte la seconde section.

Alors apparaît \(\varepsilon_R\). Non comme une décimale implorée par la
clôture, mais comme la différence que la structure était tenue de
transporter :
\[
\epsone=\frac{20}{81}\Delta_4-(S_E-1).
\]
La retenue publie ses chiffres ; la profondeur certifie sa limite ; la carte
angulaire le reconnaît comme
\[
5.13830167585752820716851646226459474899\ldots
\times10^{-8}\,{}^\circ.
\]
L'équation du radian est clôturée parce que les deux membres sont des publications du
même état :
\BigRadionEquation

L'histoire ne s'arrête pas à l'égalité. Le couple \((A,C^*)\) construit une
ellipse positive. Son produit de demi-axes fixe \(2\hret\) ; son quotient fixe la
forme. Le théorème d'aire démontre qu'une unité de phase balaie \(\hret\) et
que le tour enferme \(\hfull\). C'est ici que la division
par \(2\pi\) acquiert un contenu physique : \(h\) est l'action de tour ; \(\hbar\) est l'action par radion.

Les foyers cessent d'être des ornements d'une figure. Leur séparation reconstruit
\(C^*\) dans la carte spectrale, et leur distance de Klein mesure une profondeur
hyperbolique finie au sein d'un domaine dont l'aire hyperbolique totale est
infinie. La même frontière contient une action finie. Le paradoxe apparent
est la thèse centrale de la géométrie : une mesure clôt l'action ; une autre ouvre la
profondeur.

Le cercle n'est donc pas faux. C'est la projection exacte qui conserve la phase
et oublie l'anisotropie, l'intérieur et la mémoire. L'ellipse conserve l'action et la forme.
Klein conserve la profondeur. L'hélice conserve le retour et la biographie. Aucune
image ne remplace les autres ; toutes naissent du même état.

La réalisation de Hilbert reconnaît l'ellipse comme contour de covariance
minimale. L'oscillateur LC la reconnaît comme échange de réserves électrique et
magnétique. Le lecteur \(90/120\) publie perméabilité, permittivité,
impédance et vitesse. L'incidence six--huit produit un écran causal ;
Hodge convertit l'orientation en dualité des champs ; Lorentz convertit champ et
charge en dynamique ; Cartan--Holst reçoit l'holonomie en géométrie
gravitationnelle. L'ordre importe : l'objet HMT sélectionne ; la physique
conventionnelle réalise.

Le mot \emph{radion} se trouve ainsi justifié. Il ne signifie pas qu'un foyer soit
une charge positive et l'autre une charge négative. Il signifie que l'unité de
phase possède une orientation signée avant la publication électrique :
\[
\mathcal A(\mathfrak r_+)=+\hret,
\qquad
\mathcal A(\mathfrak r_-)=-\hret.
\]
La charge apparaît lorsque la mécanique dimensionnelle applique le transducteur. Le signe
n'est pas ajouté à la fin ; il est conservé depuis le feuillet.

On explique aussi pourquoi le secteur \(90/120\) semblait cacher de nombreuses
réponses. Incidence, extracteur dodécaphasique, queue spectrale, noyau de
Barbero, tours et retenue partagent une origine. C'est précisément parce qu'ils partagent une origine
qu'ils peuvent être comparés. C'est précisément parce qu'ils ont des opérations distinctes qu'ils ne doivent pas
être fusionnés. Le typage ne réduit pas l'unité du système : il la rend intelligible.

% REMATE_LOCAL_RADION_CIERRE
\Needspace{28\baselineskip}
La conclusion la plus compacte de cet ouvrage peut s'écrire en quatre lignes :
\[
\boxed{
\begin{aligned}
\text{phase :}\quad&\theta=1,\\
\text{action :}\quad&\mathcal A=\sigma\hret,\\
\text{mémoire :}\quad&H_9(n+9)=H_9(n)+(0,1),\\
\text{compatibilité :}\quad&1=S_E-\Phi_T+\epsone.
\end{aligned}}
\]

Le radion est cette unité quadruple. Le radian est son ombre exacte dans la carte
angulaire. La différence n'est pas dans la valeur de \(\pi\) ; elle réside dans la quantité de
structure que nous sommes disposés à conserver lorsque nous disons qu'un tour s'est
fermé.

\vspace{2\baselineskip}
\begin{center}
\begin{minipage}{.82\textwidth}
\centering\large\itshape\color{RadNavy}
Le cercle publie la phase.\\
L'ellipse contient l'action.\\
Klein ouvre la profondeur.\\
L'hélice se souvient du chemin.\\[.7em]
Le radion est l'unité qui n'oblige pas à choisir entre elles.
\end{minipage}
\end{center}

